% !TEX program = xelatex
% !TeX encoding = utf8
% !TeX spellcheck = pl-PL
%
% Praca magisterska — UniVerify
% Szablon: EE-dyplom (Wydział Elektryczny PW, aut. Łukasz Makowski, CC-BY 4.0)
%          https://github.com/SP5LMA/EE-dyplom
% Kompilacja: XeLaTeX + biber (na Overleaf ustaw kompilator na XeLaTeX).

\documentclass[thesis=mgr,faculty=ee]{EE-dyplom}

%%%%%%%%%%%%%%%%%%%%%%%%%%%%%%%%%%%%%%%%%%%%%%%%%%%%%%%%%%%%%%%%%%%%%%%%%%%
% Konfiguracja strony tytułowej — do personalizacji
%%%%%%%%%%%%%%%%%%%%%%%%%%%%%%%%%%%%%%%%%%%%%%%%%%%%%%%%%%%%%%%%%%%%%%%%%%%
\instytut{Instytut Elektrotechniki Teoretycznej i Systemów Informacyjno-Pomiarowych}
\kierunek{Informatyka Stosowana}
\specjalnosc{Inżynieria Oprogramowania} % TODO: potwierdź specjalność dla mgr
\title{Optymalizacja kosztów weryfikacji dyplomów w łańcuchu bloków Ethereum}
\engtitle{Cost Optimization of Diploma Verification on the Ethereum Blockchain}
\album{323777}
\author{inż. Nazar Shcherbyna}
\promotor{dr hab. inż. Bartosz Sawicki, prof. uczelni}
\date{2026}
\longdate{2026-XX-XX} % TODO: data złożenia

%%%%%%%%%%%%%%%%%%%%%%%%%%%%%%%%%%%%%%%%%%%%%%%%%%%%%%%%%%%%%%%%%%%%%%%%%%%
% Streszczenie (PL) i Abstract (EN)
%%%%%%%%%%%%%%%%%%%%%%%%%%%%%%%%%%%%%%%%%%%%%%%%%%%%%%%%%%%%%%%%%%%%%%%%%%%
\streszczeniepracy{
Weryfikacja autentyczności dyplomów uczelni wyższych opiera się dziś przede
wszystkim na bezpośrednim kontakcie z instytucją wydającą --- procesie powolnym,
kosztownym i uzależnionym od ciągłości jej działania. Niniejsza praca przedstawia
system UniVerify, który wykorzystuje publiczny łańcuch bloków Ethereum jako
rejestr odporny na nieuprawnioną modyfikację. Dyplomy grupowane są w~drzewa
Merkle, których korzenie publikowane są w~smart contract, dzięki czemu
koszt zapisu nie zależy od liczby dokumentów w~partii.

Praca wnosi trzy zasadnicze rozszerzenia względem systemu podstawowego. Po
pierwsze, wprowadza mechanizm selektywnego ujawniania danych oparty na
kryptograficznych zobowiązaniach z~solą, pozwalający posiadaczowi dyplomu
udostępnić jedynie wybrane pola, zgodnie z~zasadą minimalizacji danych. Po
drugie, przedstawia pomiary kosztu wdrożenia w~czterech sieciach: warstwie
pierwszej (Sepolia) oraz trzech rollupach warstwy drugiej (Arbitrum, Base,
zkSync), wraz z~analizą wrażliwości na wahania cen. Po trzecie, proponuje model
zarządzania uprawnieniami łączący portfel wielopodpisowy z~mechanizmem opóźnienia
czasowego, który eliminuje pojedynczy punkt zaufania.

Wyniki pomiarów pokazują, że koszt rejestracji pojedynczego dyplomu w~warstwie
drugiej wynosi od około $9\cdot10^{-7}$~USD w~najtańszej z~badanych sieci (Base)
do rzędu $10^{-5}$~USD w~pozostałych rollupach (partia 1000~dyplomów). Pracę zamyka analiza wartości: porównanie
proponowanego rozwiązania z~podejściem tradycyjnym oraz z~infrastrukturą klucza
publicznego, prowadzące do wniosku, że marginalna przewaga łańcucha bloków
sprowadza się do jawnego, niemodyfikowalnego i~możliwego do zweryfikowania
rejestru uprawnień do wydawania dyplomów.
}

\slowakluczowe{łańcuch bloków, Ethereum, weryfikacja dyplomów, drzewa Merkle,
selektywne ujawnianie danych, ochrona prywatności, rollupy warstwy drugiej,
smart contracts, zarządzanie uprawnieniami}

\thesisabstract{
The authentication of university diplomas today relies chiefly on direct contact
with the issuing institution --- a process that is slow, costly, and dependent on
that institution's continued operation. This thesis presents UniVerify, a system
that uses the public Ethereum blockchain as a tamper-resistant registry. Diplomas
are grouped into Merkle trees whose roots are published in a smart contract, so
that the on-chain cost is independent of the number of documents in a batch.

The thesis makes three principal extensions to the base system. First, it
introduces a selective-disclosure mechanism based on salted cryptographic
commitments, allowing a diploma holder to reveal only chosen fields, in line with
the data-minimisation principle. Second, it reports deployment-cost measurements
across four networks: layer one (Sepolia) and three layer-two rollups (Arbitrum,
Base, zkSync), together with a sensitivity analysis under price fluctuations.
Third, it proposes an authorisation-governance model that combines a
multisignature wallet with a time-lock, removing the single point of trust.

The measurements show that registering a single diploma on layer two costs from
about $9\cdot10^{-7}$~USD on the cheapest network studied (Base) to the order of
$10^{-5}$~USD on the remaining rollups (batch of 1000 diplomas). The thesis closes with a value analysis: a comparison of
the proposed solution with the status quo and with public-key infrastructure,
leading to the conclusion that the marginal advantage of the blockchain reduces to
a public, immutable, and verifiable log of diploma-issuing authorisations.
}

\thesiskeywords{blockchain, Ethereum, diploma verification, Merkle trees,
selective disclosure, privacy, layer-two rollups, smart contracts,
authorisation governance}

%%%%%%%%%%%%%%%%%%%%%%%%%%%%%%%%%%%%%%%%%%%%%%%%%%%%%%%%%%%%%%%%%%%%%%%%%%%
% Makra pomocnicze
%%%%%%%%%%%%%%%%%%%%%%%%%%%%%%%%%%%%%%%%%%%%%%%%%%%%%%%%%%%%%%%%%%%%%%%%%%%
\newcommand{\code}[1]{\texttt{#1}}

% Wykaz skrótów bez numerów stron (przy \glsaddall byłyby to martwe "1")
\renewcommand{\glossaryentrynumbers}[1]{}

% URL-e w bibliografii bez rozciągania międzyznakowego
\setlength{\biburlnumskip}{0mu}
\setlength{\biburlucskip}{0mu}
\setlength{\biburllcskip}{0mu}

% Wykresy i tabele z danych (rozdz. 4-5 — koszt gazu + benchmarki L2)
\usepackage{pgfplots}
\usepackage{pgfplotstable}
\usetikzlibrary{positioning, arrows.meta, fit, calc}
% Wspólny styl diagramów blokowych (rozdz. 3, 6, 7)
\tikzset{
  blok/.style={draw, rounded corners=2pt, align=center, font=\small,
               minimum height=8mm, inner sep=5pt},
  blokszary/.style={blok, fill=gray!10},
  blokakcent/.style={blok, fill=blue!8},
  strz/.style={-{Stealth[length=2.5mm]}, thick},
  etyk/.style={font=\footnotesize, align=center},
}
\usepackage{booktabs}
\usepackage{float}
\usepackage{multirow}   % \multirow (rozdz. 4)
\usepackage{makecell}   % \makecell (rozdz. 4)
\usepackage{etoolbox}
\pgfplotsset{compat=1.18}

% Spójny styl tabel: treść o krok mniejsza (podpisy zostają w rozmiarze tekstu,
% dla spójności z rysunkami). Węższe odstępy kolumn ograniczają przepełnienia.
\AtBeginEnvironment{tabular}{\small}
\setlength{\tabcolsep}{5pt}
\newcommand{\DataDir}{dane/l2} % katalog z danymi pomiarowymi

% Spis treści tylko do poziomu sekcji (krótszy, bez sieroty na 3. stronie)
\setcounter{tocdepth}{1}

% Bibliografia bez justowania — zapobiega rozciąganiu URL-i
\renewcommand*{\bibfont}{\normalfont\raggedright}

\begin{document}
    % Wszystkie skróty z tekst/glossary.tex trafiają do wykazu
    \glsaddall

    % Strony nagłówkowe (tytułowa + streszczenie + abstract + oświadczenie)
    \frontpages

    % Treść pracy — rozdziały w katalogu tekst/
    \chapter{Wstęp}
\label{ch:wstep}

% Cel rozdziału: wprowadzić problem, cel, zakres i wkład pracy.
% Do napisania od zera (rozdział krótki, 3-5 stron).

\section{Motywacja}
\label{sec:motywacja}
% Problem weryfikacji dyplomów: fałszerstwa, koszt weryfikacji ręcznej,
% zależność od wydawcy (uczelnia musi żyć i odpowiadać na zapytania).
\TODO{Opisz problem weryfikacji dyplomów i jego znaczenie.}

\section{Cel i zakres pracy}
\label{sec:cel}
% Cel: zaprojektować i zmierzyć tani, prywatny rejestr dyplomów na Ethereum.
% Zakres: L1 + L2, prywatność (selektywne ujawnianie), zarządzanie.
\TODO{Sprecyzuj cel i granice zakresu.}

\section{Wkład własny}
\label{sec:wklad}
% Trzy główne wkłady (mapują się na rozdziały 5-7):
\begin{enumerate}
  \item Selektywne ujawnianie pól przez zobowiązania kryptograficzne
        (rozdz.~\ref{ch:prywatnosc}).
  \item Wdrożenie i pomiary kosztów na sieciach warstwy~2
        (rozdz.~\ref{ch:l2}).
  \item Model zarządzania (multisig + timelock) i analiza wartości
        blockchaina względem alternatyw (rozdz.~\ref{ch:governance}).
\end{enumerate}
\TODO{Rozwiń każdy wkład w 1-2 zdaniach.}

\section{Struktura pracy}
\label{sec:struktura}
\TODO{Krótki przewodnik po rozdziałach 2--8.}

    \chapter{Podstawy teoretyczne}
\label{ch:podstawy}

% Jedyny rozdział czysto teoretyczny — do napisania od zera.
% Zbiera pojęcia potrzebne w rozdziałach 3-7.

\section{Weryfikacja dyplomów: stan obecny}
\label{sec:stan-obecny}
% Metody obecne: telefon/e-mail do uczelni, centralne rejestry, PKI.
\TODO{Przegląd obecnych metod weryfikacji i ich ograniczeń.}

\section{Drzewa Merkle}
\label{sec:merkle}
% Definicja, dowód przynależności (proof), złożoność O(log n),
% domenowo separowany liść.
\TODO{Formalne wprowadzenie drzew Merkle i dowodów przynależności.}

\section{Schematy zobowiązań kryptograficznych}
\label{sec:zobowiazania}
% Własności: hiding + binding. Zobowiązanie z solą: c = keccak(wartość, sól).
% Podstawa selektywnego ujawniania (rozdz. 6).
\TODO{Wprowadź zobowiązania hash-owe z solą; własności hiding/binding.}

\section{Rollupy warstwy 2 i EIP-4844}
\label{sec:l2-podstawy}
% Optimistic vs zk rollupy, dostępność danych (DA), blob transactions (blobs),
% podział kosztu: opłata L1-data + wykonanie L2.
\TODO{Model kosztowy rollupów i wpływ EIP-4844 (blobów) na koszt DA.}

\section{Infrastruktura klucza publicznego (PKI) i podpisy cyfrowe}
\label{sec:pki}
% Baza porównawcza dla analizy wartości w rozdz. 7:
% podpisany dokument daje offline'ową, niezależną weryfikację.
\TODO{Opisz PKI jako punkt odniesienia dla analizy wartości blockchaina.}

\section{Ochrona danych osobowych (RODO) i minimalizacja danych}
\label{sec:rodo}
% Zasada minimalizacji danych — motywacja dla selektywnego ujawniania.
\TODO{Zasada minimalizacji danych z RODO jako motywacja prywatności.}

    \chapter{Architektura systemu UniVerify}
\label{ch:architektura}

Rozdział przedstawia budowę systemu UniVerify w~postaci, na której opierają się
kolejne rozdziały pomiarowe i~projektowe. Opisuje model danych, kontrakt rejestru,
pakiety współdzielone oraz aplikację webową, a~na końcu zbiera je w~trzech
przepływach end-to-end: wydania, weryfikacji i~rewokacji. Pomiary kosztów pozostają
poza zakresem tego rozdziału --- omawiają je rozdziały~\ref{ch:gaz} i~\ref{ch:l2}.

\section{Model danych}
\label{sec:model-danych}

Podstawową jednostką zapisu on-chain nie jest pojedynczy dyplom, lecz \emph{partia}
(\emph{batch}). Wystawca grupuje dyplomy w~drzewo Merkle i~publikuje w~kontrakcie
jedynie jego korzeń. Dzięki temu koszt zapisu nie zależy od liczby dokumentów
w~partii --- jego uzasadnienie ilościowe podaje rozdział~\ref{ch:gaz}. Trójka
identyfikująca partię to:

\begin{itemize}
  \item \code{issuer} --- adres Ethereum portfela autoryzowanego przez uczelnię;
  \item \code{batchId} --- 64-bitowy numer kolejny partii, unikalny w~obrębie
        wystawcy;
  \item \code{merkleRoot} --- 32-bajtowy korzeń drzewa Merkle danej partii.
\end{itemize}

Pojedynczy dyplom reprezentowany jest przez \emph{kopertę}
(\emph{envelope}) złożoną z~dwóch części: \code{payload} --- danych dyplomu
(identyfikator uczelni, dane studenta, stopień, data wydania, numer dyplomu) ---
oraz \code{proof} typu \code{MERKLE\_BATCH}, zawierającego ścieżkę Merkle wiążącą
liść dokumentu z~korzeniem partii. Weryfikacja polega na odtworzeniu korzenia
z~liścia i~ścieżki oraz porównaniu go z~korzeniem odczytanym z~kontraktu.

Status dyplomu opisuje wartość \code{StatusCode}: $0$~---~nieznany
(\emph{Unknown}), $1$~---~ważny (\emph{Valid}), $2$~---~unieważniony
(\emph{Revoked}). Status nieznany oznacza, że liść nie należy do żadnej znanej
partii; unieważnienie omówiono w~sekcji~\ref{sec:przeplywy}.

\section{Kontrakt \code{DiplomaRegistry}}
\label{sec:kontrakt}

Kontrakt \code{DiplomaRegistry} udostępnia funkcje na trzech poziomach uprawnień.

\begin{description}
  \item[Właściciel (\emph{owner}).] Zarządza uczelniami i~wystawcami: rejestruje
        uczelnię wraz z~metadanymi (nazwa, kraj, strona, identyfikator
        akredytacji), przypisuje wystawcom uczelnie oraz odbiera im uprawnienia.
        Właściciel korzysta z~dwustopniowego przekazania własności
        (\code{transferOwnership} $\rightarrow$ \code{acceptOwnership}), co
        rozdział~\ref{ch:governance} wykorzystuje do zainstalowania modelu
        zarządzania bez zmiany kontraktu.
  \item[Wystawca (\emph{issuer}).] Publikuje korzenie partii
        (\code{issueBatchRoot}) oraz unieważnia pojedyncze dyplomy
        (\code{revokeFromBatch}), przy czym unieważnienie wymaga podania dowodu
        Merkle potwierdzającego przynależność dokumentu do partii.
  \item[Publiczny odczyt (\emph{view}).] Funkcje bezstanowe dostępne dla każdego:
        \code{statusWithProof} (status wraz z~weryfikacją dowodu), \code{getBatch}
        (korzeń partii), \code{isIssuer} oraz \code{issuerUniversityId}.
\end{description}

Kontrakt utrzymuje metadane uczelni, mapowanie wystawca~$\rightarrow$~uczelnia,
korzenie Merkle indeksowane parą \code{(issuer, batchId)} oraz zbiór unieważnień.
Status uczelni opisuje osobne wyliczenie: nieznana~(0), aktywna~(1),
zawieszona~(2), unieważniona~(3).

\section{Pakiety współdzielone}
\label{sec:pakiety}

System zorganizowano jako monorepozytorium, w~którym logika kryptograficzna żyje
w~jednym miejscu i~jest współdzielona przez aplikację webową, CLI oraz SDK
(rys.~\ref{fig:architektura}).

\begin{figure}[H]
\centering
\begin{tikzpicture}[node distance=6mm and 8mm]
  % Warstwa klientów
  \node[blok] (web) {aplikacja webowa\\\code{Next.js}};
  \node[blok, right=of web] (cli) {CLI\\\code{issue / verify / gov}};
  \node[blok, right=of cli] (sdk) {SDK\\\code{UniverifySdk}};
  % Rdzeń
  \node[blokakcent, below=9mm of cli, minimum width=68mm]
    (core) {\code{verifier-core}\\haszowanie, drzewa Merkle, schemat ładunku, ABI};
  % On-chain
  \node[blokszary, below=9mm of core, minimum width=68mm]
    (chain) {kontrakt \code{DiplomaRegistry} (L1 / L2)\\korzenie partii, uprawnienia wystawców, unieważnienia};
  \draw[strz] (web) -- (core);
  \draw[strz] (cli) -- (core);
  \draw[strz] (sdk) -- (core);
  \draw[strz] (core) -- node[etyk, right=1mm] {odczyt / transakcje RPC} (chain);
\end{tikzpicture}
\caption[Architektura systemu UniVerify]{Architektura systemu UniVerify: trzy interfejsy klienckie współdzielą
jedną implementację kryptograficzną (\code{verifier-core}), a~na łańcuchu
utrzymywany jest wyłącznie kontrakt rejestru.}
\label{fig:architektura}
\end{figure}

\begin{description}
  \item[\code{verifier-core}.] Warstwa fundamentalna. Eksportuje
        \code{hashPayload} (kanonizacja + Keccak256), \code{computeMerkleLeaf}
        (haszowanie liścia z~separacją domen), \code{buildMerkleFromLeaves}
        (budowa drzewa i~zwrot korzenia oraz dowodów), \code{computeRootFromProof}
        (odtworzenie korzenia z~liścia i~dowodu), a~także ABI kontraktu i~schemat
        walidacji ładunku (Zod). Musi być zbudowana przed pozostałymi pakietami.
  \item[\code{sdk}.] Klasa \code{UniverifySdk} opakowuje odczyty z~kontraktu
        i~pełną weryfikację: \code{verify} (weryfikacja koperty wraz z~metadanymi
        uczelni), \code{status}, \code{getBatch}, \code{getIssuer},
        \code{getUniversity} oraz odczyt informacji o~własności.
  \item[\code{verifier-cli}.] Skrypty wiersza poleceń: \code{issue.ts}
        i~\code{verify.ts}, a~także grupa poleceń \code{gov} obsługująca model
        zarządzania z~rozdziału~\ref{ch:governance}.
\end{description}

Rozdzielenie logiki do \code{verifier-core} sprawia, że ta sama, przetestowana
implementacja haszowania i~drzew Merkle jest używana zarówno przy wystawianiu, jak
i~przy weryfikacji --- eliminuje to ryzyko rozjazdu między stroną wydającą
a~sprawdzającą.

\section{Aplikacja webowa}
\label{sec:web}

Interfejs użytkownika to aplikacja Next.js udostępniająca trasy dopasowane do ról
w~systemie (tab.~\ref{tab:web-routes}) oraz dwa punkty API.

\begin{table}[H]
\centering
\caption{Trasy aplikacji webowej i~ich role.}
\label{tab:web-routes}
\begin{tabular}{@{}lll@{}}
\toprule
Trasa & Rola & Przeznaczenie \\
\midrule
\code{/issuer}   & wystawca  & budowa partii Merkle, publikacja korzeni, eksport dowodów \\
\code{/verifier} & publiczna & weryfikacja koperty i~sprawdzenie statusu on-chain \\
\code{/revoke}   & wystawca  & unieważnianie dyplomów przez dowód Merkle \\
\code{/present}  & posiadacz & selektywne ujawnianie pól (rozdz.~\ref{ch:prywatnosc}) \\
\code{/admin}    & właściciel & zarządzanie wystawcami i~uczelniami (rozdz.~\ref{ch:governance}) \\
\bottomrule
\end{tabular}
\end{table}

Punkty \code{/api/verify} (pełna weryfikacja z~metadanymi uczelni) oraz
\code{/api/status} (lekki odczyt statusu) chroni wspólny \emph{guard}: opcjonalne
uwierzytelnianie kluczem API (wymuszane tylko, gdy ustawiono zmienną środowiskową
z~kluczami; akceptuje nagłówek \code{Bearer} lub \code{X-API-Key}) oraz ograniczanie
liczby żądań (30~żądań na~60~s na~adres IP, w~przesuwnym oknie w~pamięci).

\section{Przepływy: wydanie, weryfikacja, rewokacja}
\label{sec:przeplywy}

Trzy główne przepływy spinają opisane elementy w~całość
(rys.~\ref{fig:przeplywy}).

\begin{figure}[H]
\centering
\begin{tikzpicture}[node distance=14mm and 24mm]
  % Wydanie
  \node[blok] (uczelnia) {uczelnia\\(wystawca)};
  \node[blokakcent, right=of uczelnia] (drzewo) {drzewo Merkle\\(off-chain, bez gazu)};
  \node[blokszary, right=of drzewo] (kontrakt) {kontrakt\\\code{DiplomaRegistry}};
  \node[blok, below=of drzewo] (absolwent) {absolwent\\(koperta: \code{payload} + dowód)};
  \node[blok, below=of kontrakt] (weryfikator) {weryfikator};
  \draw[strz] (uczelnia) -- node[etyk, above=1mm] {liście\\dyplomów} (drzewo);
  \draw[strz] (drzewo) -- node[etyk, above=1mm] {\code{issueBatchRoot}\\(korzeń)} (kontrakt);
  \draw[strz] (drzewo) -- node[etyk, left=1mm] {eksport kopert} (absolwent);
  \draw[strz] (absolwent) -- node[etyk, above=1mm] {koperta} (weryfikator);
  \draw[strz] (weryfikator) -- node[etyk, right=1mm, align=left]
    {\code{statusWithProof}:\\porównanie korzeni} (kontrakt);
  \draw[strz, dashed] (uczelnia.south) -- ++(0,-32mm) -|
    node[etyk, below, pos=0.25] {\code{revokeFromBatch} (dowód Merkle)}
    (kontrakt.south east);
\end{tikzpicture}
\caption[Przepływy: wydanie, weryfikacja, rewokacja]{Przepływy w~systemie: wydanie (uczelnia publikuje wyłącznie korzeń
partii), weryfikacja (odtworzenie korzenia z~koperty i~porównanie ze stanem
kontraktu) oraz rewokacja (linia przerywana).}
\label{fig:przeplywy}
\end{figure}

\paragraph{Wydanie.} Wystawca buduje off-chain drzewo Merkle z~liści dyplomów
(\code{computeMerkleLeaf} + \code{buildMerkleFromLeaves}), publikuje w~kontrakcie
sam korzeń (\code{issueBatchRoot}) i~eksportuje każdemu absolwentowi kopertę zawierającą
jego \code{payload} oraz dowód Merkle. Cała praca kryptograficzna poza jedną
transakcją odbywa się bez gazu.

\paragraph{Weryfikacja.} Weryfikator (przez \code{/verifier}, API lub SDK)
odtwarza z~koperty liść, z~liścia i~dowodu --- korzeń, po czym porównuje go
z~korzeniem odczytanym z~kontraktu dla pary \code{(issuer, batchId)}. Pozytywny
wynik oznacza, że dokument należy do partii opublikowanej przez autoryzowanego
wystawcę i~nie został unieważniony.

\paragraph{Rewokacja.} Aby unieważnić pojedynczy dyplom, wystawca wywołuje
\code{revokeFromBatch}, podając dowód Merkle dokumentu. Kontrakt dodaje liść do
zbioru unieważnień, przez co kolejne odczyty statusu zwracają \code{Revoked}, mimo
że korzeń partii pozostaje niezmieniony. Pozwala to cofnąć pojedynczy dokument bez
naruszania pozostałych dyplomów tej samej partii.

    \chapter{Optymalizacja kosztu gazu w warstwie pierwszej}
\label{ch:gaz}

Rozdział bada, jak decyzje projektowe wpływają na koszt gazu kontraktu rejestru
dyplomów na Ethereum. Eksperymenty podzielono na trzy etapy, każdy odpowiadający
innemu poziomowi abstrakcji decyzji: (1)~mikrooptymalizacje modelu per-dyplom ---
siedem kolejnych wersji kontraktu z~pojedynczymi zmianami kodu; (2)~zmianę
architektury na \emph{Merkle batch}, w~której jedna transakcja pokrywa całą partię;
(3)~mikrooptymalizacje hot-path weryfikacji dowodu Merkle. Wnioski z~trzech
etapów zbiera sekcja~\ref{sec:gaz-wnioski}.

\section{Metodyka pomiaru}
\label{sec:gaz-metodyka}

Wszystkie pomiary wykonano lokalnie w~środowisku Foundry~\cite{foundry}
(Solidity 0.8.20), co eliminuje szum wynikający ze zmienności opłat sieciowych. Foundry dostarcza dwie
granulacje danych:

\begin{itemize}
  \item \textbf{Suite-level} --- łączny gaz jednego testu, włącznie z~setupem
        (deploy kontraktu, wywołania przygotowawcze). Oddaje realistyczny koszt
        całego scenariusza.
  \item \textbf{Function-level} (\code{--gas-report}) --- gaz samej funkcji EVM,
        agregowany po wszystkich jej wywołaniach w~suite jako min/avg/max. Pozwala
        precyzyjnie izolować koszt jednej operacji.
\end{itemize}

\noindent Komendy użyte do pomiarów:
\begin{verbatim}
cd contracts
forge test --match-path "test/DiplomaRegistry[ABCDEF]*.t.sol" --gas-report
forge test --match-path "test/DiplomaRegistryGasComparison.t.sol" --gas-report
\end{verbatim}

\section{Etap 1 --- mikrooptymalizacje modelu per-dyplom}
\label{sec:gaz-etap1}

Model per-dyplom to najprostsze możliwe podejście: każdy dyplom to jeden wpis
\code{mapping(bytes32 => Record)}. Zbadano siedem kolejnych wersji kontraktu, aby
sprawdzić, jak bardzo zmiana jednej konkretnej decyzji implementacyjnej wpływa na
koszt gazu. Wersjom nadano nazwy opisujące \emph{kluczową właściwość}
odróżniającą każdą z~nich od poprzedniej.

\subsection{Wersje kontraktu}

\begin{table}[H]
\centering
\caption{Chronologiczne wersje kontraktu per-dyplom: zakres i~uzasadnienie zmian.}
\begin{tabular}{>{\bfseries}p{3.0cm}
                >{\raggedright\arraybackslash}p{4.5cm}
                >{\raggedright\arraybackslash}p{6.5cm}}
\toprule
Nazwa wersji & Kluczowa zmiana & Uzasadnienie / hipoteza \\
\midrule
\code{StringErrors}
  & \code{require(..., "STRING")}, \newline rekord: \code{\{issuer, issuedAt, revoked\}}, \newline mutable \code{owner}
  & Punkt startowy --- kontrakt pisany intuicyjnie, bez optymalizacji \\[6pt]

\code{TimestampRevoke}
  & \code{revokedAt(uint64)} zamiast \code{bool}
  & Hipoteza: timestamp rewokacji przyda się audytorom; \newline \textbf{efekt: regres} --- dodatkowy slot storage \\[6pt]

\code{CustomErrors}
  & \code{custom errors} zamiast komunikatów string
  & Komunikaty string są kodowane jako dane, \newline custom errors to 4-bajtowy selektor \\[6pt]

\code{ImmutableOwner}
  & \code{owner} jako \code{immutable}
  & \code{immutable} trzyma wartość w~bytecode, \newline nie w~storage --- brak odczytu SLOAD w~modyfikatorze \\[6pt]

\code{Idempotent}
  & \code{addIssuer}/\code{removeIssuer} nie rzucają przy ponownym wywołaniu
  & Zabezpieczenie przed kosztownym revert-em \newline w~przypadku race condition \\[6pt]

\code{MinimalRecord}
  & Usunięcie timestampów ze storage; \newline rekord: \code{\{issuer, revoked\}}
  & Daty pozostają dostępne w~logach zdarzeń; \newline duplikacja w~storage jest zbędna \\[6pt]

\code{PackedStorage}
  & \code{mapping(bytes32 => uint256)}: \newline issuer (160 bitów) + bit revoked \newline w~jednym slocie; dodano \code{status()}
  & Jeden SLOAD zamiast dwóch --- rozkład adresu \newline i~flagi z~jednej 256-bitowej wartości \\
\bottomrule
\end{tabular}
\end{table}

\subsection{Wyniki}

Poniższa tabela pokazuje łączny koszt jednego scenariusza testowego (suite-level).
Warto pamiętać, że \code{testGas\_revoke} zawiera wywołanie \code{issue} jako setup
stanu --- dlatego wartości w~tej kolumnie są wyższe niż sam koszt rewokacji.

\begin{table}[H]
\centering
\caption{Koszt testów end-to-end w~gazie, łącznie z~deployem i~setupem (kolumny to
warianty \code{testGas\_addIssuer}, \code{testGas\_issue}, \code{testGas\_revoke}).}
\begin{tabular}{l r r r >{\raggedright\arraybackslash}p{3.1cm}}
\toprule
Wersja & \code{addIssuer} & \code{issue} & \code{revoke} & Uwagi \\
\midrule
\code{StringErrors}    & 57\,119 & 60\,721 &  93\,116 & punkt startowy \\
\code{TimestampRevoke} & 57\,153 & 62\,890 & 114\,348 & \textbf{regres: +21k gas na revoke} \\
\code{CustomErrors}    & 57\,186 & 63\,012 & 114\,380 & pozornie bez poprawy (por.\ function-level) \\
\code{ImmutableOwner}  & 55\,100 & 63\,037 & 114\,431 & drobna poprawa addIssuer \\
\code{Idempotent}      & 55\,312 & 63\,047 & 114\,448 & marginalnie droższy w~tym scenariuszu \\
\code{MinimalRecord}   & 55\,100 & 60\,569 &  92\,868 & powrót do taniego \code{revoke} \\
\code{PackedStorage}   & 55\,056 & 60\,531 &  92\,797 & najlepszy wynik ogólny \\
\bottomrule
\end{tabular}
\end{table}

\noindent Najistotniejszą obserwacją jest \textbf{regres w~wersji
\code{TimestampRevoke}}: dodanie pola \code{uint64 revokedAt} do rekordu zwiększa
koszt rewokacji z~93k do 114k gas, ponieważ EVM musi zaktualizować dodatkowy slot
storage przy każdym \code{revoke()}. Wersja \code{MinimalRecord} naprawia ten błąd,
usuwając timestampy ze storage i~pozostawiając je wyłącznie w~logach zdarzeń.

Widok function-level pozwala zmierzyć czysty koszt każdej operacji, niezależnie od
kontekstu testu. Różnice deployment są tutaj bardziej widoczne niż w~suite-level.

\begin{table}[H]
\centering
\caption{Koszt deploymentu i~średni koszt operacji wg \code{forge --gas-report} (Avg).}
\begin{tabular}{lrrrr}
\toprule
Wersja & Deployment & \code{addIssuer} & \code{issue} & \code{revoke} \\
\midrule
\code{StringErrors}    & 464\,211 & 45\,719 & 48\,458 & 31\,274 \\
\code{TimestampRevoke} & 469\,181 & 46\,963 & 50\,659 & 50\,369 \\
\code{CustomErrors}    & 409\,497 & 45\,904 & 43\,367 & 38\,867 \\
\code{ImmutableOwner}  & 403\,109 & 43\,366 & 43\,370 & 38\,870 \\
\code{Idempotent}      & 418\,029 & 43\,523 & 43\,377 & 38\,874 \\
\code{MinimalRecord}   & 373\,484 & 43\,366 & 41\,607 & 28\,825 \\
\code{PackedStorage}   & 393\,596 & 43\,621 & 43\,084 & 29\,302 \\
\bottomrule
\end{tabular}
\end{table}

\noindent Widok function-level ujawnia efekt niewidoczny na poziomie suite:
\code{CustomErrors}
redukuje deployment z~464k do 409k gas, ponieważ komunikaty string (\code{"Only owner"},
\code{"Already issued"}, itd.) są przechowywane w~bytecode kontraktu, a~\code{custom
errors} zastępuje je 4-bajtowymi selektorami. Ta zmiana oszczędza \emph{rozmiar
kodu}, co przekłada się na tańszy deployment: redukcja o~55\,000 gas (11{,}9\%).

Nie każda zmiana poprawia sytuację --- \code{TimestampRevoke} jest przykładem
modyfikacji wprowadzonej w~dobrej intencji, której koszt znacząco przewyższył
oczekiwania: przechowywanie znacznika czasu rewokacji w~storage zamiast
w~zdarzeniach zwiększyło koszt \code{revoke} o~61\%. Sumarycznie, od \code{StringErrors} do \code{MinimalRecord}:
\[
  \Delta_{\code{issue}} \approx 14{,}1\%,\quad
  \Delta_{\code{revoke}} \approx 7{,}8\%,\quad
  \Delta_{\text{deploy}} \approx 19{,}6\%.
\]

\section{Etap 2 --- architektura Merkle batch}
\label{sec:gaz-etap2}

Model per-dyplom zakłada, że \emph{każdy dyplom trafia na łańcuch jako osobna
transakcja}. Architektura Merkle batch znosi to założenie: zamiast zapisywać $N$
dyplomów w~$N$ transakcjach, wystawca:
\begin{enumerate}
  \item buduje drzewo Merkle z~$N$ liści (off-chain, bez gazu),
  \item zapisuje na łańcuch wyłącznie \textbf{korzeń drzewa} --- jedną transakcję
        \code{issueBatchRoot(batchId, merkleRoot)},
  \item weryfikacja dyplomu przez \code{statusWithProof(docHash, issuer, batchId,
        proof)} sprawdza przynależność liścia do korzenia, używając dowodu
        $O(\log N)$.
\end{enumerate}

\subsection{Zmiana modelu operacji}

\begin{table}[H]
\centering
\caption{Mapowanie operacji: model per-dyplom (\code{PackedStorage}) vs Merkle.}
\begin{tabular}{p{3.5cm} p{4.5cm} p{5.5cm}}
\toprule
Operacja & \code{PackedStorage} (per-dyplom) & \code{MerkleBaseline} \\
\midrule
Wystawienie & \code{issue(docHash)} --- 1 tx / dyplom & \code{issueBatchRoot(batchId, root)} --- 1 tx / N dyplomów \\[4pt]
Weryfikacja & \code{status(docHash)} --- odczyt mapping & \code{statusWithProof(hash, issuer, batchId, proof)} --- weryfikacja Merkle \\[4pt]
Rewokacja & \code{revoke(docHash)} --- prosty SSTORE & \code{revokeFromBatch(hash, batchId, proof)} --- SSTORE + weryfikacja \\[4pt]
Zarządzanie \newline wystawcami & \code{addIssuer}/\code{removeIssuer} & \code{onboardIssuerAndUniversity} (bogatszy model z~uczelnią) \\
\bottomrule
\end{tabular}
\end{table}

\subsection{Wyniki: koszt na dyplom}

\begin{table}[H]
\centering
\caption{\code{PackedStorage} vs \code{MerkleBaseline} --- min/avg/max wg \code{forge --gas-report}.}
\begin{tabular}{llrrr}
\toprule
Model & Operacja & Min & Avg & Max \\
\midrule
\multirow{4}{*}{\code{PackedStorage}} &
  \code{addIssuer}    & 21\,585 & 43\,621 & 44\,855 \\
& \code{issue}        & 24\,030 & 43\,084 & 48\,200 \\
& \code{removeIssuer} & 21\,607 & 23\,112 & 24\,947 \\
& \code{revoke}       & 26\,407 & 29\,302 & 31\,222 \\
\midrule
\multirow{4}{*}{\code{MerkleBaseline}} &
  \code{issueBatchRoot}  & 50\,927 & 50\,937 & 50\,939 \\
& \code{revokeFromBatch} & 54\,239 & 57\,229 & 61\,591 \\
& \code{statusWithProof} & 10\,555 & 11\,890 & 13\,835 \\
& \code{onboardIssuer\&Uni} & \multicolumn{3}{c}{76\,907 (stałe)} \\
\midrule
\multicolumn{2}{l}{Deployment} & \multicolumn{2}{l}{\code{PackedStorage}: 393\,596} & \\
\multicolumn{2}{l}{} & \multicolumn{2}{l}{\code{MerkleBaseline}: 1\,059\,364 ($\times$2{,}7)} & \\
\bottomrule
\end{tabular}
\end{table}

\noindent Na poziomie pojedynczej transakcji Merkle nie jest tańszy ---
\code{issueBatchRoot} kosztuje tyle samo co \code{issue} w~per-dyplom. Właściwą
metryką porównania jest jednak \textbf{koszt na dyplom}, ponieważ jedna
transakcja \code{issueBatchRoot} obejmuje $N$ dyplomów:
\[
  \text{gas\_per\_diploma}(N) = \frac{50\,939}{N},
  \qquad
  \text{gas\_per\_diploma}_\text{legacy} = 48\,200 = \text{const}.
\]

\begin{figure}[H]
\centering
\begin{tikzpicture}
\begin{axis}[
    width=0.90\textwidth,
    height=8cm,
    xlabel={Rozmiar partii $N$},
    ylabel={Gas na dyplom (skala logarytmiczna)},
    xmode=log, ymode=log,
    log basis x={10}, log basis y={10},
    xmin=1, xmax=1000,
    ymin=10, ymax=100000,
    grid=both,
    minor grid style={gray!15},
    major grid style={gray!40},
    legend style={at={(0.5,-0.18)},anchor=north,font=\small},
    legend columns=2,
]
\addplot+[mark=*, thick, color=blue!70!black, mark size=3pt]
  coordinates {(1,50939)(2,25469)(5,10188)(10,5094)(20,2547)
               (50,1019)(100,509)(200,255)(500,102)(1000,51)};
\addlegendentry{\code{MerkleBaseline}: $50\,939/N$}
\addplot+[mark=square*, thick, color=orange!85!black, dashed, mark size=3pt]
  coordinates {(1,48200)(2,48200)(5,48200)(10,48200)(20,48200)
               (50,48200)(100,48200)(200,48200)(500,48200)(1000,48200)};
\addlegendentry{\code{PackedStorage}: stałe 48\,200 gas}
\draw[dashed,gray] (axis cs:100,1) -- (axis cs:100,509);
\node[anchor=south west, font=\footnotesize] at (axis cs:100,509)
  {$N=100$: $509$ vs $48\,200$};
\end{axis}
\end{tikzpicture}
\caption[Amortyzacja kosztu publikacji: legacy vs Merkle batch]{Amortyzacja kosztu publikacji: model Merkle obniża koszt na dyplom
odwrotnie proporcjonalnie do rozmiaru partii $N$, podczas gdy w~modelu
per-dyplom koszt jest stały; próg opłacalności leży już przy $N \approx 1$
(skala logarytmiczna na obu osiach).}
\end{figure}

\begin{table}[H]
\centering
\caption{Konkretne wartości: gas na dyplom i~stosunek do modelu legacy.}
\begin{tabular}{rrrr}
\toprule
$N$ & Merkle (gas/dyplom) & Legacy (gas/dyplom) & Krotność redukcji \\
\midrule
1    & 50\,939 & 48\,200 & $0{,}95\times$ (legacy tańszy) \\
5    & 10\,188 & 48\,200 & $4{,}7\times$ \\
10   &  5\,094 & 48\,200 & $9{,}5\times$ \\
50   &  1\,019 & 48\,200 & $47\times$ \\
100  &    509  & 48\,200 & $\mathbf{94{,}7\times}$ \\
500  &    102  & 48\,200 & $472\times$ \\
1000 &     51  & 48\,200 & $945\times$ \\
\bottomrule
\end{tabular}
\end{table}

\noindent Przejście na Merkle to \emph{nie kolejna mikrooptymalizacja} --- to inna
odpowiedź na pytanie o~model danych. Koszt jednostkowej transakcji jest taki sam jak
legacy (lub nawet minimalnie wyższy), ale pojedyncza transakcja ,,certyfikuje'' całą
partię. Przy $N=100$ koszt per dyplom spada z~$\approx$48\,200 do $\approx$509 gas
--- \textbf{redukcja o 98{,}9\%}. Ceną za tę zmianę jest niemal trzykrotnie droższy
deployment ($\approx$1{,}06M vs $\approx$0{,}39M gas) oraz bardziej złożona weryfikacja
wymagająca dowodu Merkle.

\section{Etap 3 --- optymalizacje hot-path Merkle}
\label{sec:gaz-etap3}

Po wyborze architektury Merkle pozostaje pytanie, czy ścieżkę weryfikacji dowodu
można dalej zoptymalizować na poziomie kodu EVM. Kontrakt
\code{DiplomaRegistryG} (dalej: \code{MerkleG}) to identyczny funkcjonalnie wariant
\code{MerkleBaseline} z~czterema przyrostowymi technikami nałożonymi na dwie
wewnętrzne funkcje: \code{\_processProof()} (iteracja po elementach dowodu)
i~\code{\_merkleLeaf()} (budowanie liścia do zhashowania).

\subsection{Techniki G1--G4}

Tabela~\ref{tab:g1g4} zestawia cztery zastosowane techniki wraz z~miejscem ich
działania i~mechanizmem oszczędności.

\begin{table}[H]
\centering
\caption{Cztery techniki zastosowane w~\code{MerkleG}.}
\label{tab:g1g4}
\begin{tabular}{>{\bfseries}p{2.8cm} p{4.0cm} p{7.0cm}}
\toprule
Nazwa & Gdzie & Co robi i~dlaczego oszczędza gas \\
\midrule
G1 \newline \code{unchecked++i}
  & \code{\_processProof}, pętla
  & Blok \code{unchecked\{\}} eliminuje sprawdzanie przepełnienia licznika
    iteracji; pre-increment \code{++i} zamiast \code{i++} nie tworzy kopii. Razem:
    kilka opcodów mniej na krok. \\[4pt]

G2 \newline assembly keccak
  & \code{\_processProof}, krok haszowania
  & Hashowanie pary liści przez \code{keccak256(0x00, 0x40)} --- dane pisane
    bezpośrednio do \emph{scratch space} EVM (bajty 0x00--0x3f), który istnieje
    zawsze i~nie wymaga aktualizacji wskaźnika wolnej pamięci
    (\code{mload/mstore 0x40}). Eliminuje alokację pamięci na każdym kroku pętli. \\[4pt]

G3 \newline prywatny liść
  & \code{\_merkleLeaf}, ścieżka wewnętrzna
  & Publiczny \code{merkleLeaf()} weryfikuje, że \code{docHash != 0} i~\code{issuer
    != 0}. Prywatne \code{\_merkleLeaf()} te sprawdzenia pomija --- bo kod
    wywołujący je wcześniej już zweryfikował dane. Oszczędność: dwa warunkowe rzuty
    mniej. \\[4pt]

G4 \newline assembly \mbox{preimage}
  & \code{\_merkleLeaf}, budowanie danych
  & Zamiast \code{abi.encodePacked(address(this), chainid(), issuer, batchId,
    docHash)} (112 bajtów, alokacja przez ABI-encoder), dane pakowane są przez
    ręczne \code{mstore} z~przesunięciami bitowymi. Oszczędność: cały narzut
    ABI-encodera na budowanie 112-bajtowego bufora. \\
\bottomrule
\end{tabular}
\end{table}

\subsection{Wyniki}

\begin{table}[H]
\centering
\caption{\code{MerkleBaseline} vs \code{MerkleG} --- min/avg/max wg Foundry.}
\begin{tabular}{lrrrrrrc}
\toprule
Operacja
  & \multicolumn{3}{c}{\code{MerkleBaseline}}
  & \multicolumn{3}{c}{\code{MerkleG}}
  & $\Delta$ Avg \\
\cmidrule(lr){2-4}\cmidrule(lr){5-7}
& Min & Avg & Max & Min & Avg & Max & \\
\midrule
Deployment
  & \multicolumn{3}{c}{1\,059\,364}
  & \multicolumn{3}{c}{1\,018\,273}
  & $-41\,091$ (\textbf{$-$3{,}9\%}) \\
\code{issueBatchRoot}
  & 50\,927 & 50\,937 & 50\,939
  & 50\,927 & 50\,937 & 50\,939
  & $0$ \\
\code{revokeFromBatch}
  & 54\,239 & 57\,229 & 61\,591
  & 54\,101 & 56\,443 & 59\,859
  & $-786$ ($-1{,}4\%$) \\
\code{statusWithProof}
  & 10\,555 & 11\,890 & 13\,835
  & 10\,417 & 11\,104 & 12\,103
  & $-786$ (\textbf{$-$6{,}6\%}) \\
\code{merkleLeaf} (view)
  & 866 & 866 & 866
  & 824 & 824 & 824
  & $-42$ ($-4{,}8\%$) \\
\bottomrule
\end{tabular}
\end{table}

\noindent \code{issueBatchRoot} nie wywołuje \code{\_processProof} ani
\code{\_merkleLeaf} w~ścieżce zapisu, więc oszczędność tam wynosi zero.
\code{statusWithProof} zyskuje relatywnie więcej ($-6{,}6\%$ avg) niż
\code{revokeFromBatch} ($-1{,}4\%$ avg), ponieważ stanowi mniejszą bazę kosztową (mniej
operacji storage) --- absolutna oszczędność na wywołaniu wynosi jednak \textbf{786
gas avg} dla obu operacji.

\paragraph{Co dokładnie mierzy test.}
Przed analizą liczb w~funkcji głębokości drzewa warto zrozumieć, co naprawdę
robi każdy test. Na przykładzie \code{testRevokeFromBatch\_Depth8}:

\begin{enumerate}
  \item \textbf{\code{\_makeLeaves(256)}} --- funkcja testowa wywołuje
        \code{reg.merkleLeaf()} \emph{256 razy} jako zewnętrzne wywołania kontraktu,
        żeby uzyskać zahaszowane liście. Każde kosztuje 866 gas. Łącznie:
        $256 \times 866 = 221\,696$ gas.
  \item \textbf{\code{\_buildTree(leaves)}} --- buduje drzewo Merkle w~czystym
        Solidity (funkcja \code{pure}): oblicza 255 wewnętrznych węzłów używając
        \code{keccak256(abi.encodePacked(l, r))} w~zagnieżdżonej pętli. Solidity
        alokuje tymczasowy bufor w~pamięci EVM na każde wywołanie
        \code{abi.encodePacked}, a~dostęp do dwuwymiarowej tablicy
        \code{bytes32[][]} wymaga podwójnego sprawdzenia granic przy każdym
        odczycie. Efekt: $\approx$9\,000 gas na węzeł, czyli
        $255 \times 9\,000 \approx 2\,200\,000$ gas łącznie.
  \item \textbf{\code{reg.issueBatchRoot(BATCH, root)}} --- 50\,939 gas.
  \item \textbf{\code{reg.revokeFromBatch(doc, BATCH, proof)}} z~dowodem
        8-elementowym --- 61\,591 gas (baseline) / 59\,859 gas (G).
\end{enumerate}

\noindent\textbf{Kluczowa obserwacja:} kroki 1 i~2 to \emph{artefakt testu}.
Symulują one coś, co w~produkcji odbywa się \textbf{off-chain, bez żadnego gazu}.
Uczelnia buduje drzewo Merkle lokalnie (np.\ w~TypeScript/Python) i~jedynie zapisuje
jego korzeń na łańcuch --- jedną transakcją \code{issueBatchRoot}. Dlatego liczby
z~tabeli poniżej \textbf{nie} odpowiadają kosztowi produkcyjnemu.

\begin{table}[H]
\centering
\caption{Rozkład kosztu \code{testRevokeFromBatch\_DepthD} na składowe (Baseline).}
\begin{tabular}{crrrrrr}
\toprule
$d$ & Liści & \makecell{Budowa liści\\$2^d \times 866$} & \makecell{Budowa drzewa\\w Solidity} & \makecell{\code{issueBatch}\\Root} & \makecell{\code{revoke}\\FromBatch} & \textbf{Suma} \\
\midrule
0 & 1   &    866 &         0 & 50\,939 & 54\,239 & 124\,856\phantom{*} \\
1 & 2   &  1\,732 &    11\,077 & 50\,939 & 55\,158 & 137\,905\phantom{*} \\
4 & 16  & 13\,856 &   134\,697 & 50\,939 & 57\,915 & 276\,384\phantom{*} \\
8 & 256 & 221\,696 & 2\,219\,313 & 50\,939 & 61\,591 & 2\,553\,539\phantom{*} \\
\bottomrule
\end{tabular}
\end{table}

\noindent Koszty \code{revokeFromBatch} dla głębokości 1 i~4 są interpolowane na
podstawie zmierzonego przyrostu 919 gas/krok z~danych function-level. Suma dla
każdego wiersza zgadza się ze zmierzoną wartością testu. W~wierszu $d=8$ budowa
drzewa ($\approx$2{,}2M gas) stanowi \textbf{87\% łącznego kosztu testu}, podczas gdy
rzeczywiste wywołanie \code{revokeFromBatch} kosztuje jedynie 61\,591 gas
(\textbf{2{,}4\% całości}).

Rozziew między kosztem testu a~kosztem produkcyjnym rośnie z~głębokością: przy
$d=8$ test (2\,553\,539 gas) jest niemal 23-krotnie droższy od rzeczywistej
pary operacji on-chain (\code{issueBatchRoot} + \code{revokeFromBatch},
112\,530 gas) --- cała ta różnica to koszt budowania drzewa Merkle w~Solidity,
co w~produkcji odbywa się off-chain.

\paragraph{Wniosek praktyczny.} Niezależnie od tego, czy partia liczy 1~czy
256~dyplomów, koszt produkcyjny jednej rewokacji to stałe
$\approx$105--113~tys.\ gas (\code{issueBatchRoot} + \code{revokeFromBatch}).
Rozmiar drzewa zmienia koszt testu, ale nie zmienia znacząco kosztu on-chain.

\noindent Dla kompletności, poniższa tabela zestawia pełne wyniki testów
end-to-end dla obu operacji i~obu wersji kontraktu:

\begin{table}[H]
\centering
\caption{Pełne wyniki testów end-to-end --- obie operacje, obie wersje.}
\begin{tabular}{crrrr}
\toprule
Głębokość $d$ & Liści $2^d$ & \code{MerkleBaseline} & \code{MerkleG} & Oszczędność \\
\midrule
\multicolumn{5}{l}{\itshape revokeFromBatch (koszt testu, $\gg$ koszt produkcyjny przy $d \geq 4$)} \\
\midrule
0 & 1   & 124\,856    & 124\,676    & 180 \\
1 & 2   & 137\,905    & 137\,484    & 421 \\
4 & 16  & 276\,384    & 274\,778    & 1\,606 \\
8 & 256 & 2\,553\,539 & 2\,541\,055 & 12\,484 \\
\midrule
\multicolumn{5}{l}{\itshape statusWithProof (koszt testu)} \\
\midrule
0 & 1   & 81\,106     & 80\,926     & 180 \\
1 & 2   & 93\,558     & 93\,137     & 421 \\
4 & 16  & 230\,521    & 228\,915    & 1\,606 \\
8 & 256 & 2\,505\,681 & 2\,493\,197 & 12\,484 \\
\bottomrule
\end{tabular}
\end{table}

\subsection{Model liniowy oszczędności}

Obserwując jak oszczędność rośnie z~głębokością, można ją rozłożyć na trzy składniki:

\begin{equation}
\Delta(d) =
  \underbrace{\bigl(2^d + 1\bigr) \cdot 42}_{\text{wywołania \code{merkleLeaf} w~teście (G4)}}
  + \underbrace{96}_{\text{ominięte zero-checks (G3)}}
  + \underbrace{d \cdot 199}_{\text{kroki pętli dowodu (G1+G2)}},
  \label{eq:model}
\end{equation}

gdzie:
\begin{itemize}
  \item $42$ gas --- oszczędność jednego publicznego wywołania \code{merkleLeaf}
        (assembly zamiast \code{abi.encodePacked}, G4); test wykonuje $2^d$ takich
        wywołań przy budowie liści oraz jedno dodatkowe dla dokumentu badanego;
  \item $96$ gas --- oszczędność na wewnętrznym wywołaniu \code{\_merkleLeaf}
        (G4 + pominięte sprawdzenia G3);
  \item $199$ gas --- oszczędność na jednym kroku pętli \code{\_processProof} (G1+G2).
\end{itemize}

\noindent\textbf{Weryfikacja modelu}:

\begin{table}[H]
\centering
\caption{Model \eqref{eq:model} vs zmierzone wartości.}
\begin{tabular}{crrr}
\toprule
$d$ & $\Delta(d)$ (model) & Zmierzone & Różnica \\
\midrule
0 & $2\cdot42 + 96 + 0\cdot199 = 180$    & 180    & $\mathbf{0}$ \\
1 & $3\cdot42 + 96 + 1\cdot199 = 421$    & 421    & $\mathbf{0}$ \\
4 & $17\cdot42 + 96 + 4\cdot199 = 1\,606$  & 1\,606 & $\mathbf{0}$ \\
8 & $257\cdot42 + 96 + 8\cdot199 = 12\,482$ & 12\,484 & $+2$ \\
\bottomrule
\end{tabular}
\end{table}

\noindent Model \eqref{eq:model} odtwarza wyniki pomiarów dokładnie dla
$d \in \{0, 1, 4\}$; przy $d=8$ pozostaje resztkowe $2$~gas (poniżej
$0{,}02\%$ wartości), przypisywalne nieliniowemu kosztowi rozszerzania pamięci
EVM przy 256~liściach. Potwierdza to, że oszczędności wynikają z~trzech
niezależnych i~przewidywalnych źródeł, a~nie z~niezamierzonych efektów
ubocznych.

\section{Koszt w~przeliczeniu na PLN}
\label{sec:gaz-pln}

Przyjęto kurs \textbf{1 ETH = 7000 PLN}, spójny z~przyjętym
w~rozdziale~\ref{ch:l2} kursem $\text{ETH}/\text{USD}=3500$ (migawka
z~13~maja~2026~r.) przy kursie $\text{USD}/\text{PLN}\approx2{,}00$, oraz trzy
scenariusze ceny gazu (niski: 20 gwei; typowy: 40 gwei; szczytowy: 80 gwei).
Wzór przeliczenia:
\[
\text{PLN} = \text{gas} \times (\text{gasPrice}_{[\text{gwei}]} \times 10^{-9})
             \times \text{kurs}_{[\text{PLN/ETH}]}.
\]

\begin{table}[H]
\centering
\caption{Koszt kluczowych operacji w~PLN (kurs: 7000 PLN/ETH).}
\begin{tabular}{lrrrr}
\toprule
Operacja & Gas (max) & 20 gwei & 40 gwei & 80 gwei \\
\midrule
\code{PackedStorage issue}          & 48\,200  & 6{,}75 & 13{,}50 & 26{,}99 \\
\code{MerkleBaseline issueBatchRoot}& 50\,939  & 7{,}13 & 14{,}26 & 28{,}53 \\
\code{MerkleG revokeFromBatch}      & 59\,859  & 8{,}38 & 16{,}76 & 33{,}52 \\
\code{MerkleG statusWithProof}      & 12\,103  & 1{,}69 & 3{,}39 & 6{,}78 \\
Deployment \code{MerkleG}           & 1\,018\,273 & 142{,}56 & 285{,}12 & 570{,}23 \\
\bottomrule
\end{tabular}
\end{table}

\noindent\textbf{Przykład amortyzacji} przy 40 gwei i~partii $N=100$: jeden
\code{issueBatchRoot} kosztuje ${\approx}14{,}26$~PLN, co po podziale na
100~dyplomów daje ${\approx}0{,}14$~PLN na dyplom --- wobec $13{,}50$~PLN na
dyplom w~modelu per-dyplom. Wystawienie jednego dyplomu jest więc przy typowej
cenie gazu i~partii 100~absolwentów niemal stukrotnie tańsze.

\section{Wnioski etapu}
\label{sec:gaz-wnioski}

Trzy etapy eksperymentów układają się w~spójny obraz: decyzje projektowe mają różne
\emph{rzędy wpływu} na koszt.

\begin{table}[H]
\centering
\caption{Zestawienie wyników trzech etapów.}
\begin{tabular}{>{\raggedright\arraybackslash}p{4.0cm} >{\raggedright\arraybackslash}p{3.4cm} r >{\raggedright\arraybackslash}p{3.2cm}}
\toprule
Etap & Typ zmiany & Redukcja kosztu & Metryka \\
\midrule
1: \code{StringErrors} $\to$ \code{MinimalRecord} & mikrooptymalizacje kodu   & $\approx$14\%     & gas/dyplom (avg) \\
1: deployment                                    & mikrooptymalizacje kodu   & $\approx$20\%     & deployment gas \\
2: \code{PackedStorage} $\to$ \code{MerkleBaseline} & zmiana architektury      & $\approx$98{,}9\%   & gas/dyplom przy $N=100$ \\
3: \code{MerkleBaseline} $\to$ \code{MerkleG}       & mikrooptymalizacje hot-path & $\approx$6{,}6\%  & \code{statusWithProof} avg \\
\bottomrule
\end{tabular}
\end{table}

Cztery wnioski wynikają z~całości. Po pierwsze, \textbf{mikrooptymalizacje kodu
dają realną, ale ograniczoną poprawę}: różnice rzędu 8--20\% są mierzalne i~warte
uwagi przy wdrożeniu na głównej sieci Ethereum, ale nie zmieniają klasy kosztowej
--- każdy dyplom nadal kosztuje dziesiątki tysięcy gas. Po drugie, \textbf{zmiana
architektury działa o~rząd wielkości silniej}: model Merkle batch nie poprawia
istniejącego kodu, lecz zmienia model danych na poziomie projektu. Przy partii 100 dyplomów koszt
jednostkowy spada o~niemal 99\%, co jest efektem amortyzacji stałego kosztu
transakcji, a~nie szybszego algorytmu. Po trzecie, \textbf{optymalizacje hot-path
Merkle są przewidywalne i~liniowe}: wzór $\Delta(d) = (2^d+1) \cdot 42 + 96 + d \cdot
199$ opisuje zmierzone oszczędności z~dokładnością do 2~gas --- wiemy \emph{skąd} pochodzi każdy
zaoszczędzony gas i~możemy ekstrapolować na inne rozmiary drzew. Po czwarte,
\textbf{dobór optymalizacji zależy od modelu użycia}: dla systemu z~małymi partiami
($N < 5$) różnice architektur są marginalne, dla masowej emisji ($N \geq 50$)
Merkle jest jedynym sensownym wyborem, a~optymalizacje G1--G4 mają sens przy
systemach z~intensywną weryfikacją lub bardzo głębokimi drzewami.

Pomiary mają jednak dwa ograniczenia: są lokalne (Foundry) i~nie uwzględniają
rzeczywistej zmienności \emph{base fee} ani kosztów calldata (które rosną
z~długością dowodu) --- rozdział~\ref{ch:l2} uzupełnia je o~pomiary sieciowe na
warstwie drugiej; ponadto nie mierzono kosztu budowania drzewa Merkle po stronie
klienta (off-chain), co przy bardzo dużych batchach ($N > 10^5$) może być istotnym
parametrem.

    \chapter{Wdrożenie i pomiary kosztów w sieciach warstwy drugiej}
\label{ch:l2}

Rozdział~\ref{ch:gaz} wykazał, że po stronie samego kontraktu koszt publikacji
korzenia Merkle jest niemal stały i sprowadza się do pojedynczej, taniej operacji
zapisu. Rzeczywisty koszt ponoszony przez uczelnię zależy jednak nie tylko od
zużycia gazu, lecz przede wszystkim od sieci, w~której kontrakt zostaje wdrożony.
Warstwa pierwsza Ethereum zapewnia najsilniejsze gwarancje bezpieczeństwa, ale jej
koszt bywa o~rzędy wielkości wyższy niż w~sieciach warstwy drugiej (ang.
\emph{layer~2}, L2), które przenoszą wykonanie poza łańcuch główny, publikując na
nim jedynie skompresowane dane. Niniejszy rozdział mierzy ten koszt empirycznie:
wdraża identyczny kontrakt \code{DiplomaRegistry} w~czterech sieciach i porównuje
koszt kluczowych operacji zarówno w~jednostkach gazu, jak i w~dolarach
amerykańskich.

\section{Metodyka pomiaru}
\label{sec:l2-metodyka}

\subsection{Zakres i badane sieci}
Porównanie obejmuje koszt uruchomienia kontraktu \code{DiplomaRegistry} na czterech
sieciach Ethereum: warstwie pierwszej (Sepolia, jako punkt odniesienia) oraz trzech
rozwiązaniach warstwy drugiej reprezentujących odmienne rodziny technologiczne ---
Arbitrum Nitro, Base jako przedstawiciel stosu OP oraz zkSync~Era jako rollup
oparty na dowodach zwięzłości (ang. \emph{validity proofs}). Dobór ten pozwala
porównać nie pojedyncze wdrożenia, lecz reprezentatywne klasy rollupów.

\subsection{Pomiary}
Każdą sieć mierzono na jej publicznej sieci testowej. \emph{Zużycie gazu} oraz
\emph{opłata za dane L1} są wielkościami rzeczywistymi, odczytanymi z~paragonów
transakcji, natomiast koszty w~USD modelowane są na podstawie historycznych
wartości \code{baseFeePerGas} warstwy pierwszej z~ostatnich 90~dni (percentyle
p10/p50/p90). Wszystkie wyniki w~USD przeliczono kursem
$\text{ETH}/\text{USD}=3500$ (snapshot z~dnia 13~maja~2026~r.); wpływ tego założenia
badany jest w~sekcji~\ref{sec:l2-wrazliwosc}. Pomiary wykonano 19~maja~2026~r.

\subsection{Parametry próby}
Dla operacji \code{issueBatch} użyto rozmiarów paczek
$\{1, 10, 100, 1\,000, 10\,000\}$. Dla \code{revokeFromBatch} zebrano 30~próbek
przeciwko paczce bazowej o~rozmiarze 40, z~losowym wyborem unieważnianych liści.
Dla \code{statusWithProof} zebrano 100~próbek czasu odpowiedzi RPC. Każdy pomiar
opóźnienia powtórzono $N{=}30$ razy, a~raportowane wartości są średnią z~trzech
niezależnych przebiegów pomiarowych.

\section{Koszt operacji \code{issueBatch}}
\label{sec:l2-issue}

Operacja \code{issueBatch} publikuje korzeń Merkle całej paczki dyplomów.
Tabela~\ref{tab:issue-gas} pokazuje jej zużycie gazu, a~tabela~\ref{tab:issue-usd}
--- odpowiadający mu koszt w~USD. Ponieważ na łańcuchu zapisywany jest zawsze
pojedynczy 32-bajtowy korzeń, zużycie gazu jest praktycznie niezależne od liczby
dyplomów w~paczce; różnicę kosztu między sieciami tworzy sposób rozliczania danych,
nie rozmiar danych.

\begin{table}[H]
  \centering
  \caption{Zużycie gazu operacji \code{issueBatch} według sieci i~rozmiaru paczki
  (wartości deterministyczne).}
  \label{tab:issue-gas}
  \pgfplotstabletypeset[
    col sep=comma, string type,
    every head row/.style={before row=\toprule, after row=\midrule},
    every last row/.style={after row=\bottomrule},
    columns={chain,gas_1,gas_10,gas_100,gas_1000,gas_10000},
    columns/chain/.style={column name=Sieć, column type=l},
    columns/gas_1/.style={column name={$n{=}1$}, column type=r},
    columns/gas_10/.style={column name={$n{=}10$}, column type=r},
    columns/gas_100/.style={column name={$n{=}100$}, column type=r},
    columns/gas_1000/.style={column name={$n{=}1000$}, column type=r},
    columns/gas_10000/.style={column name={$n{=}10000$}, column type=r},
  ]{\DataDir/table-issue-cost.csv}
\end{table}

\begin{table}[H]
  \centering
  \caption{Koszt operacji \code{issueBatch} w~USD (scenariusz p50,
  $\text{ETH}/\text{USD}{=}3500$). Pełną zmienność scenariuszy p10/p50/p90 dla
  $n{=}1000$ przedstawia rys.~\ref{fig:basefee}.}
  \label{tab:issue-usd}
  \pgfplotstabletypeset[
    col sep=comma, string type,
    every head row/.style={before row=\toprule, after row=\midrule},
    every last row/.style={after row=\bottomrule},
    columns={chain,usd_p50_1,usd_p50_10,usd_p50_100,usd_p50_1000,usd_p50_10000},
    columns/chain/.style={column name=Sieć, column type=l},
    columns/usd_p50_1/.style={column name={$n{=}1$}, column type=r},
    columns/usd_p50_10/.style={column name={$n{=}10$}, column type=r},
    columns/usd_p50_100/.style={column name={$n{=}100$}, column type=r},
    columns/usd_p50_1000/.style={column name={$n{=}1000$}, column type=r},
    columns/usd_p50_10000/.style={column name={$n{=}10000$}, column type=r},
  ]{\DataDir/table-issue-cost.csv}
\end{table}

Praktyczną miarą jest koszt w~przeliczeniu na jeden dyplom, malejący wraz ze
wzrostem paczki dzięki amortyzacji stałego kosztu publikacji korzenia
(tabela~\ref{tab:per-diploma}, rys.~\ref{fig:cost-per-diploma}).

\begin{table}[H]
  \centering
  \caption{Koszt na dyplom (amortyzacja paczki) w~USD, scenariusz p50.}
  \label{tab:per-diploma}
  \pgfplotstabletypeset[
    col sep=comma, string type,
    every head row/.style={before row=\toprule, after row=\midrule},
    every last row/.style={after row=\bottomrule},
    columns/chain/.style={column name=Sieć, column type=l},
    columns/usd_p50_per_diploma_1/.style={column name={$n{=}1$}, column type=r},
    columns/usd_p50_per_diploma_10/.style={column name={$n{=}10$}, column type=r},
    columns/usd_p50_per_diploma_100/.style={column name={$n{=}100$}, column type=r},
    columns/usd_p50_per_diploma_1000/.style={column name={$n{=}1000$}, column type=r},
    columns/usd_p50_per_diploma_10000/.style={column name={$n{=}10000$}, column type=r},
  ]{\DataDir/table-issue-per-diploma.csv}
\end{table}

\begin{figure}[H]
  \centering
  \begin{tikzpicture}
    \begin{loglogaxis}[
      xlabel={Rozmiar paczki (liczba dyplomów)},
      ylabel={Koszt na dyplom (USD)},
      legend pos=north east, grid=major,
      width=0.85\textwidth, height=8cm,
    ]
      \addplot table[x=batchSize, y=sepolia] {\DataDir/fig1-cost-per-diploma.dat};
      \addplot table[x=batchSize, y=arbitrumSepolia] {\DataDir/fig1-cost-per-diploma.dat};
      \addplot table[x=batchSize, y=baseSepolia] {\DataDir/fig1-cost-per-diploma.dat};
      \addplot table[x=batchSize, y=zksyncSepolia] {\DataDir/fig1-cost-per-diploma.dat};
      \legend{Sepolia L1, Arbitrum Sepolia, Base Sepolia, zkSync Sepolia}
    \end{loglogaxis}
  \end{tikzpicture}
  \caption{Koszt na dyplom w~funkcji rozmiaru paczki (skala log-log, scenariusz
  p50).}
  \label{fig:cost-per-diploma}
\end{figure}

\section{Koszt operacji \code{revokeFromBatch}}
\label{sec:l2-revoke}

Unieważnienie pojedynczego dyplomu (\code{revokeFromBatch}) jest operacją rzadką i
dotyczy jednej transakcji. Tabela~\ref{tab:revoke} podaje jej medianowy koszt;
utrzymuje on ten sam porządek sieci co koszt publikacji, z~Base jako
zdecydowanie najtańszą opcją.

\begin{table}[H]
  \centering
  \caption{Koszt operacji \code{revokeFromBatch} (pojedyncza transakcja, mediana).}
  \label{tab:revoke}
  \pgfplotstabletypeset[
    col sep=comma, string type,
    every head row/.style={before row=\toprule, after row=\midrule},
    every last row/.style={after row=\bottomrule},
    columns={chain,median_gas,median_l1data_wei,usd_p50},
    columns/chain/.style={column name=Sieć, column type=l},
    columns/median_gas/.style={column name={Gaz (mediana)}, column type=r},
    columns/median_l1data_wei/.style={column name={Dane L1 [wei]}, column type=r},
    columns/usd_p50/.style={column name={Koszt [USD]}, column type=r},
  ]{\DataDir/table-revoke-cost.csv}
\end{table}

\section{Latencja}
\label{sec:l2-latencja}

Poza kosztem istotny jest czas odpowiedzi. Tabela~\ref{tab:read-latency} podaje
latencję odczytu \code{statusWithProof} (weryfikacja po stronie sprawdzającego),
a~tabela~\ref{tab:inclusion-latency} --- czas od wysłania transakcji zapisu do
odebrania paragonu przez klienta.

\begin{table}[H]
  \centering
  \caption{Latencja odczytu \code{statusWithProof} (ms).}
  \label{tab:read-latency}
  \pgfplotstabletypeset[
    col sep=comma,
    every head row/.style={before row=\toprule, after row=\midrule},
    every last row/.style={after row=\bottomrule},
    columns={chain,p50_ms,p95_ms,max_observed_ms},
    columns/chain/.style={string type, column name=Sieć, column type=l},
    columns/p50_ms/.style={column name={$p_{50}$}, column type=r, fixed, fixed zerofill, precision=1},
    columns/p95_ms/.style={column name={$p_{95}$}, column type=r, fixed, fixed zerofill, precision=1},
    columns/max_observed_ms/.style={column name={maks.}, column type=r, fixed, fixed zerofill, precision=1},
  ]{\DataDir/table-read-latency.csv}
\end{table}

\begin{table}[H]
  \centering
  \caption{Czas odbioru paragonu transakcji zapisu po stronie klienta (ms).
  Kolumny \code{issue} pochodzą z~$N{=}30$ powtórzeń \code{issueBatch(1000)}
  w~każdym przebiegu i~są agregowane po wszystkich trzech przebiegach.}
  \label{tab:inclusion-latency}
  \pgfplotstabletypeset[
    col sep=comma, string type,
    every head row/.style={before row=\toprule, after row=\midrule},
    every last row/.style={after row=\bottomrule},
    columns/chain/.style={column name=Sieć, column type=l},
    columns/revoke_p50_ms/.style={column name={revoke $p_{50}$}, column type=r},
    columns/revoke_p95_ms/.style={column name={revoke $p_{95}$}, column type=r},
    columns/issue_p50_ms/.style={column name={issue $p_{50}$}, column type=r},
    columns/issue_p95_ms/.style={column name={issue $p_{95}$}, column type=r},
    columns/issue_sigma_ms/.style={column name={issue $\sigma$}, column type=r},
    columns/issue_n/.style={column name={$N$}, column type=r},
  ]{\DataDir/table-inclusion-latency.csv}
\end{table}

Mierzoną wielkością jest czas między wywołaniem \code{sendTransaction} a~otrzymaniem
paragonu przez klienta RPC. Przy interwale odpytywania około 4~s (domyślnym
w~bibliotece viem) wynik jest górnym ograniczeniem faktycznego czasu inkluzji
w~bloku, nie zaś jego dokładnym pomiarem; rozkład wartości na sieci Sepolia jest
wyraźnie skupiony wokół wielokrotności czasu bloku ($\approx 12$~s), co potwierdza
tę interpretację.

\section{Analiza wrażliwości}
\label{sec:l2-wrazliwosc}

Ponieważ koszty w~USD zależą od trzech zmiennych rynkowych --- opłaty bazowej L1,
kursu ETH/USD oraz ceny gazu sekwensera L2 --- poniżej badana jest odporność
wyników na ich wahania.

\subsection{Scenariusze opłaty bazowej L1}
\begin{figure}[H]
  \centering
  \begin{tikzpicture}
    \begin{axis}[
      ybar,
      symbolic x coords={sepolia, arbitrumSepolia, baseSepolia, zksyncSepolia},
      xtick=data,
      ylabel={Koszt \code{issueBatch(1000)} (USD)},
      bar width=10pt, legend pos=north west,
      width=0.85\textwidth, height=8cm,
    ]
      \addplot table[x=chain, y=quiet_usd] {\DataDir/fig4-basefee-scenarios.dat};
      \addplot table[x=chain, y=normal_usd] {\DataDir/fig4-basefee-scenarios.dat};
      \addplot table[x=chain, y=congested_usd] {\DataDir/fig4-basefee-scenarios.dat};
      \legend{p10 (spokój), p50 (typowo), p90 (kongestia)}
    \end{axis}
  \end{tikzpicture}
  \caption{Koszt \code{issueBatch(1000)} przy scenariuszach opłaty bazowej L1
  p10/p50/p90.}
  \label{fig:basefee}
\end{figure}

\subsection{Wrażliwość na kurs ETH/USD}
Koszt w~USD jest liniową funkcją kursu ETH/USD: $C(r) = C_{3500}\cdot r/3500$, gdzie
$C_{3500}$ to wartość bazowa raportowana przy snapshocie z~dnia 13~maja~2026~r.
($r=3500$). Tabela~\ref{tab:eth-sens} pokazuje koszt \code{issueBatch(1000)}
w~scenariuszu p50 dla trzech zakotwiczonych poziomów rynkowych.

\begin{table}[H]
  \centering
  \caption{Wrażliwość kosztu \code{issueBatch(1000)} na kurs ETH/USD
  (opłata bazowa p50).}
  \label{tab:eth-sens}
  \pgfplotstabletypeset[
    col sep=comma, string type,
    every head row/.style={before row=\toprule, after row=\midrule},
    every last row/.style={after row=\bottomrule},
    columns={chain,usd_2000,usd_3500,usd_5000},
    columns/chain/.style={column name=Sieć, column type=l},
    columns/usd_2000/.style={column name={2000 USD}, column type=r},
    columns/usd_3500/.style={column name={3500 USD}, column type=r},
    columns/usd_5000/.style={column name={5000 USD}, column type=r},
  ]{\DataDir/table-sensitivity-eth-usd.csv}
\end{table}

\subsection{Wrażliwość na cenę gazu sekwensera L2}
Cena gazu ustalana przez operatora sekwensera L2 jest niezależna od opłaty bazowej
L1. Model kosztu rozkłada się addytywnie:
\[
  C_{\text{total}} = C_{\text{exec}} + C_{\text{L1 data}},
\]
gdzie $C_{\text{exec}} = \code{gasUsed}\cdot p_{\text{L2}}$ skaluje się liniowo
z~ceną sekwensera $p_{\text{L2}}$, natomiast $C_{\text{L1 data}}$ pochodzi z~opłaty
blob na L1 i~pozostaje stała. Tabela~\ref{tab:l2price-sens} pokazuje koszt
\code{issueBatch(1000)} przy trzech mnożnikach względem snapshotu cen (Arbitrum
0{,}1~gwei, Base 0{,}005~gwei, zkSync 0{,}05~gwei). Sepolię pominięto --- jako sieć
L1 nie ma ceny sekwensera.

\begin{table}[H]
  \centering
  \caption{Wrażliwość kosztu \code{issueBatch(1000)} na cenę gazu sekwensera L2
  (mnożniki względem snapshotu cen).}
  \label{tab:l2price-sens}
  \pgfplotstabletypeset[
    col sep=comma, string type,
    every head row/.style={before row=\toprule, after row=\midrule},
    every last row/.style={after row=\bottomrule},
    columns={chain,usd_05x,usd_10x,usd_20x},
    columns/chain/.style={column name=Sieć, column type=l},
    columns/usd_05x/.style={column name={$0{,}5\times$}, column type=r},
    columns/usd_10x/.style={column name={$1\times$}, column type=r},
    columns/usd_20x/.style={column name={$2\times$}, column type=r},
  ]{\DataDir/table-sensitivity-l2-price.csv}
\end{table}

\section{Dyskusja}
\label{sec:l2-dyskusja}

{\sloppy
Wszystkie pomiary opisane w~niniejszym rozdziale wykonano przez jednego dostawcę
RPC w~jednej lokalizacji, z~$N{=}30$ powtórzeniami dla każdego pomiaru opóźnienia
oraz losowym wyborem liści przy unieważnianiu. Raportowane wyniki to średnia
$\pm\sigma$ z~trzech niezależnych przebiegów. Wszystkie parametry konfiguracyjne
(ziarno losowości, dostawca, region, kurs ETH/USD) zapisano w~pliku
\code{meta.json} towarzyszącym eksportowi danych.\par}

\subsection{Ograniczenia pomiaru}
Poniżej zebrano ograniczenia metodologiczne, których czytelnik powinien być
świadomy, interpretując wyniki. Cztery z~pierwotnie zidentyfikowanych ograniczeń
wyeliminowano w~toku rewizji eksperymentu (jednolity dostawca RPC, $N{=}30$
powtórzeń, losowy wybór liści do unieważnienia, trzy niezależne przebiegi
z~raportowaniem $\sigma$); pozostałe sześć omówiono poniżej.

\paragraph{Założenie kompresji Brotli dla Arbitrum.} Model kosztu zapisu danych na
L1 dla Arbitrum przyjmuje współczynnik kompresji Brotli równy 1{,}0, ponieważ dane
Merkle (hasze pseudolosowe) nie kompresują się znacząco. Realny calldata
o~strukturze rzadkiej kompresuje się 2--4-krotnie, co oznacza, że nasze
oszacowanie kosztu L1 dla Arbitrum jest \textbf{górnym ograniczeniem}; rzeczywiste
koszty produkcyjne mogą być proporcjonalnie niższe.

\paragraph{Niewspółmierność jednostki gazu w~zkSync Era.} Wartość \code{gasUsed}
raportowana przez zkSync (131\,463 dla \code{issueBatch} z~$n{=}1$) łączy w~jednej
jednostce koszt wykonania, dane publiczne oraz narzut bootloadera kontekstu Account
Abstraction --- nie jest to ta sama jednostka co \code{gasUsed} w~EVM. Porównywanie
surowych liczb gazu między rodzinami rollupów jest \textbf{metodologicznie
niepoprawne}; sensowne porównanie kosztów możliwe jest wyłącznie w~USD lub ETH.

\paragraph{Opóźnienie inkluzji ograniczone interwałem odpytywania.} Pomiar czasu od
wysłania do paragonu opiera się o~cykl odpytywania klienta viem
(\textasciitilde 4~s). Otrzymane wartości są zatem \textbf{górnym ograniczeniem
czasu inkluzji w~bloku}, a~nie samym czasem inkluzji. Porównanie międzysieciowe jest
dodatkowo zniekształcone interakcją czasu bloku z~interwałem odpytywania.

\paragraph{Snapshot skalarów Base SystemConfig.} Model kosztu Base (Ecotone) używa
snapshotu parametrów \code{baseFeeScalar} oraz \code{blobBaseFeeScalar} pobranego
z~konkretnego bloku L1. Po zmianie tych skalarów przez operatora Base wynikowe
koszty zmieniłyby się proporcjonalnie; data snapshotu jest udokumentowana
w~\code{meta.json} każdego eksportu.

\paragraph{Snapshot kursu ETH/USD.} Wyniki w~USD przeliczono kursem
$\text{ETH}/\text{USD}=3500$ (snapshot z~dnia 13~maja~2026~r.). Wynik dla innego
kursu otrzymuje się mnożeniem przez czynnik $r/3500$.

\paragraph{Konserwatywne ceny sekwensera L2.} Ceny gazu sekwensera ustawiono na
zaokrąglone wartości powyżej żywego snapshotu (Arbitrum 0{,}1 wobec 0{,}02~gwei;
Base 0{,}005 wobec 0{,}006; zkSync 0{,}05 wobec 0{,}0453), aby pochłonąć typowe
wahania kongestii.

\subsection{Reprodukowalność}
Wszystkie wyniki w~tym rozdziale można odtworzyć z~czystego klona repozytorium
UniVerify poleceniem \code{npm run benchmarks:all -{}- -{}-runs=3}, po wypełnieniu
pliku \code{.env} kluczem prywatnym portfela testowego oraz adresami RPC czterech
sieci testowych. Pipeline wykonuje sekwencję \code{deploy} $\to$ \code{measure}
$\to$ \code{price} $\to$ \code{aggregate} $\to$ \code{export}, regenerując wszystkie
tabele i~wykresy. Każdy artefakt powiązany jest z~plikiem \code{meta.json}
zawierającym identyfikator przebiegu, datę pomiaru, kurs ETH/USD wraz z~datą
snapshotu, ziarno losowości, dostawcę i~region RPC oraz listę sieci.

\subsection{Wnioski}
\label{sec:l2-wnioski}
Zebrane pomiary pozwalają odpowiedzieć na pytanie postawione na wstępie rozdziału:
przeniesienie rejestru do warstwy drugiej istotnie obniża koszt, lecz skala
korzyści zależy od rodziny rollupa. Przy realistycznym rozmiarze paczki $n{=}1000$
(rzędu rocznego wolumenu dyplomów dużego wydziału) koszt zarejestrowania jednego
dyplomu wynosi w~scenariuszu p50 około $3{,}2\cdot10^{-5}$~USD na~Sepolii (L1),
$1{,}8\cdot10^{-5}$~USD na~Arbitrum, $2{,}3\cdot10^{-5}$~USD na~zkSync oraz zaledwie
$9{,}0\cdot10^{-7}$~USD na~Base --- czyli około 35-krotnie taniej niż w~warstwie
pierwszej (por. tab.~\ref{tab:per-diploma}, rys.~\ref{fig:cost-per-diploma}).

Ponieważ koszt na dyplom maleje proporcjonalnie do $1/n$ (publikujemy jeden
32-bajtowy korzeń niezależnie od rozmiaru paczki), stosunki kosztów między sieciami
są niezmiennicze względem $n$; powyższa hierarchia utrzymuje się dla dowolnego
rozmiaru paczki. Przewaga Base nie wynika z~lepszej kompresji \emph{naszych}
danych --- we wszystkich sieciach publikujemy ten sam korzeń o~koszcie około
51~tys.\ gazu --- lecz z~tego, że pojedyncza transakcja L2 jest tania: rollup
amortyzuje koszt zapisu danych na~L1 pomiędzy wszystkich swoich użytkowników poprzez
transakcje blob (EIP-4844). Dla naszej transakcji składnik danych L1 stanowi znikomą
część kosztu, w~którym dominuje wykonanie po stronie sekwensera.

Wynik ten jest odporny na typową zmienność rynkową: analiza wrażliwości
(sekcja~\ref{sec:l2-wrazliwosc}) pokazuje, że nawet dwukrotne wahania kursu ETH/USD,
opłaty bazowej L1 czy ceny sekwensera L2 nie zmieniają uporządkowania sieci ani
rzędu wielkości kosztu. Z~perspektywy kosztowej Base jest zatem najkorzystniejszym
celem wdrożenia rejestru dyplomów spośród badanych sieci.

    \chapter{Ochrona prywatności przez selektywne ujawnianie danych}
\label{ch:prywatnosc}

System podstawowy zakotwicza w~łańcuchu bloków skrót całej treści dyplomu.
Rozdział ten pokazuje, jak --- bez żadnej zmiany w~kontrakcie --- zastąpić ten
monolityczny skrót strukturą zobowiązań kryptograficznych, dzięki której posiadacz
dyplomu może ujawnić weryfikatorowi jedynie wybrane pola, zachowując pozostałe
w~tajemnicy.

\section{Problem: minimalizacja danych}
\label{sec:priv-problem}

W~systemie podstawowym skrót dokumentu powstaje przez kanonizację i~zahaszowanie
\emph{całego} ładunku dyplomu: \code{docHash = keccak256(canonicalize(payload))}.
Aby zweryfikować dyplom, posiadacz musi przekazać weryfikatorowi pełny ładunek ---
inaczej weryfikator nie odtworzy skrótu. Oznacza to, że dowód posiadania stopnia
naukowego wymusza ujawnienie wszystkich danych: imienia i~nazwiska, numeru dyplomu,
daty wydania i~szczegółów kierunku, nawet gdy relacja weryfikacyjna wymaga
potwierdzenia tylko jednego z~nich.

Narusza to zasadę minimalizacji danych --- fundament nowoczesnych regulacji ochrony
danych osobowych --- zgodnie z~którą przetwarzać należy jedynie dane niezbędne do
celu. Pracodawca weryfikujący sam fakt ukończenia studiów nie powinien być zmuszany
do poznania numeru dyplomu kandydata. Celem tego rozdziału jest umożliwienie
weryfikacji \emph{pojedynczych pól} bez ujawniania reszty, przy zachowaniu tej samej
gwarancji zakotwiczenia w~łańcuchu.

\section{Konstrukcja zobowiązań pól}
\label{sec:priv-zobowiazania}

Ładunek dzielimy na cztery ujawnialne pola: \code{student}, \code{degree},
\code{issuedAt} oraz \code{diplomaNumber}. Dla każdego z~nich obliczamy osobne
\emph{zobowiązanie ukrywające} (\emph{hiding commitment}), łącząc kanoniczną postać
wartości pola z~losową solą:
\[
  c_i = \operatorname{keccak256}\bigl(\operatorname{canonicalize}(\text{pole}_i)
        \,\Vert\, \text{sól}_i\bigr),
  \qquad i \in \{\text{student}, \text{degree}, \text{issuedAt},
        \text{diplomaNumber}\}.
\]
Każda sól to niezależna, losowa wartość 256-bitowa generowana przez
\code{crypto.getRandomValues}. Skrót dokumentu odtwarzamy następnie \emph{ze
zobowiązań}, a~nie z~jawnych wartości:
\[
  \code{docHash} = \operatorname{keccak256}
    \bigl(c_{\text{student}} \Vert c_{\text{degree}}
    \Vert c_{\text{issuedAt}} \Vert c_{\text{diplomaNumber}}\bigr).
\]
Liść drzewa Merkle pozostaje związany z~kontekstem wdrożenia dokładnie tak, jak
w~systemie podstawowym:
\[
  \code{leaf} = \operatorname{keccak256}
    \bigl(\code{registry} \Vert \code{chainId} \Vert \code{issuer}
    \Vert \code{batchId} \Vert \code{docHash}\bigr).
\]
Ponieważ \code{docHash} jest deterministyczną funkcją czterech zobowiązań, korzeń
partii publikowany w~kontrakcie nie zmienia formatu --- kontrakt nie odróżnia
dyplomu prywatnego od standardowego.

\section{Koperta prywatna}
\label{sec:priv-koperta}

Prywatna koperta dyplomu (\code{PrivateDiplomaEnvelope}) zawiera trzy części:

\begin{itemize}
  \item \textbf{zobowiązania} --- pełny komplet czterech wartości $c_i$, potrzebny
        do odtworzenia \code{docHash} niezależnie od tego, które pola ujawniono;
  \item \textbf{pola ujawnione} --- dla każdego \emph{wybranego} pola jego jawna
        wartość wraz z~odpowiadającą jej solą; pola nieujawnione są w~kopercie
        nieobecne;
  \item \textbf{dowód Merkle} --- ścieżka od liścia do korzenia partii, identyczna
        jak w~kopercie standardowej.
\end{itemize}

Strukturę koperty i~drogę od pól do korzenia partii ilustruje
rys.~\ref{fig:koperta}.

\begin{figure}[H]
\centering
\begin{tikzpicture}[node distance=4mm and 7mm]
  \node[blok] (p1) {\code{student}\\\footnotesize + sól$_1$};
  \node[blok, right=of p1] (p2) {\code{degree}\\\footnotesize + sól$_2$};
  \node[blok, right=of p2] (p3) {\code{issuedAt}\\\footnotesize + sól$_3$};
  \node[blok, right=of p3] (p4) {\code{diplomaNumber}\\\footnotesize + sól$_4$};
  \node[blokakcent, below=7mm of p1] (c1) {$c_1$};
  \node[blokakcent, below=7mm of p2] (c2) {$c_2$};
  \node[blokakcent, below=7mm of p3] (c3) {$c_3$};
  \node[blokakcent, below=7mm of p4] (c4) {$c_4$};
  \node[blokszary, below=8mm of $(c2)!0.5!(c3)$, minimum width=50mm]
    (dochash) {\code{docHash} $= H(c_1 \Vert c_2 \Vert c_3 \Vert c_4)$};
  \node[blokszary, below=5mm of dochash, minimum width=50mm]
    (leaf) {liść $= H(\code{registry} \Vert \code{chainId} \Vert \code{issuer}
            \Vert \code{batchId} \Vert \code{docHash})$};
  \node[blokszary, below=5mm of leaf] (root) {korzeń partii (on-chain)};
  \foreach \i in {1,2,3,4}
    \draw[strz] (p\i) -- node[etyk, right=0.5mm] {$H$} (c\i);
  \draw[strz] (c1.south) -- (dochash.north west);
  \draw[strz] (c2) -- (dochash);
  \draw[strz] (c3) -- (dochash);
  \draw[strz] (c4.south) -- (dochash.north east);
  \draw[strz] (dochash) -- (leaf);
  \draw[strz] (leaf) -- node[etyk, right=0.5mm] {dowód Merkle} (root);
\end{tikzpicture}
\caption[Struktura koperty prywatnej]{Koperta prywatna: każde pole wraz z~tajną solą daje zobowiązanie
$c_i$; skrót dokumentu powstaje z~zobowiązań, nie z~jawnych wartości. Pole
ujawnione = wartość + sól; pole ukryte = samo $c_i$.}
\label{fig:koperta}
\end{figure}

Sól pełni tu podwójną rolę. Dla pola \emph{ujawnionego} jest przekazywana
weryfikatorowi, aby mógł odtworzyć i~sprawdzić zobowiązanie. Dla pola
\emph{ukrytego} pozostaje tajemnicą posiadacza --- to ona uniemożliwia odgadnięcie
wartości pola z~samego zobowiązania.

\section{Protokół weryfikacji selektywnej}
\label{sec:priv-protokol}

Weryfikator dysponujący prywatną kopertą wykonuje trzy kroki:

\begin{enumerate}
  \item \textbf{Sprawdzenie ujawnionych pól.} Dla każdego pola obecnego w~sekcji
        ujawnionej weryfikator przelicza $c_i' = \operatorname{keccak256}
        (\operatorname{canonicalize}(\text{wartość}) \Vert \text{sól})$ i~porównuje
        z~zobowiązaniem $c_i$ z~koperty. Niezgodność któregokolwiek pola oznacza
        odrzucenie.
  \item \textbf{Odtworzenie skrótu dokumentu.} Ze wszystkich czterech zobowiązań
        (ujawnionych i~ukrytych) weryfikator liczy \code{docHash}, a~z~niego liść.
  \item \textbf{Weryfikacja dowodu Merkle.} Odtworzony liść wraz z~dowodem daje
        korzeń, który weryfikator porównuje z~korzeniem odczytanym z~kontraktu dla
        pary \code{(issuer, batchId)}.
\end{enumerate}

Pole ukryte przechodzi przez krok~2 jako nieprzejrzyste zobowiązanie: jego wartość
nie jest weryfikatorowi znana, ale jego wkład do \code{docHash} --- a~więc i~do
liścia --- jest w~pełni uwzględniony. Weryfikator zyskuje pewność, że ujawnione pola
należą do dyplomu zakotwiczonego on-chain, nie dowiadując się niczego o~polach
ukrytych.

\section{Implementacja}
\label{sec:priv-implementacja}

Cała logika prywatności jest \emph{poza łańcuchem}. Kontrakt \code{DiplomaRegistry}
pozostaje niezmieniony --- publikuje korzenie przez ten sam \code{issueBatchRoot},
a~jego koszty gazu są identyczne jak w~rozdziale~\ref{ch:l2}. Zmiany ograniczają się
do warstwy klienckiej: pakiet \code{verifier-core} zyskuje funkcje
\code{generateFieldSalts}, \code{computeFieldCommitments}, \code{computePrivateDocHash}
i~\code{verifyDisclosedFields} (pokryte 17~testami jednostkowymi), SDK udostępnia
weryfikację selektywną, a~aplikacja webowa --- trasę \code{/present}, na której
posiadacz zaznacza pola do ujawnienia i~generuje prywatną kopertę.

Fakt, że tak istotna funkcja realizowana jest bez ani jednej zmiany w~kontrakcie,
jest sam w~sobie wynikiem: pokazuje, że model zakotwiczenia korzenia Merkle jest na
tyle ogólny, iż selektywne ujawnianie mieści się w~nim jako czysto klienckie
rozszerzenie.

\section{Analiza bezpieczeństwa}
\label{sec:priv-bezpieczenstwo}

Bezpieczeństwo konstrukcji opiera się na dwóch własnościach funkcji
\code{keccak256} traktowanej jako zobowiązanie.

\paragraph{Wiązanie (binding).} Odporność \code{keccak256} na kolizje sprawia, że
posiadacz nie może otworzyć jednego zobowiązania na dwie różne wartości pola.
Uniemożliwia to przedstawienie tego samego dyplomu z~podmienioną, lecz „poprawnie
weryfikującą się'' zawartością.

\paragraph{Ukrywanie (hiding).} Dla pola ukrytego napastnik widzi jedynie
$c_i = \operatorname{keccak256}(\text{wartość} \Vert \text{sól})$. Atak wyczerpujący
wymagałby zgadnięcia \emph{zarówno} wartości, \emph{jak i} 256-bitowej soli. Ponieważ
sól pochodzi z~kryptograficznie bezpiecznego generatora i~pozostaje tajna, przestrzeń
poszukiwań jest rzędu $2^{256}$ --- praktycznie nieprzeszukiwalna, nawet gdy sama
wartość pola należy do małej domeny (np. \code{degree.level} przyjmuje kilka
wartości). Odsala to typowy zarzut wobec zobowiązań nad małymi domenami: bez
tajnej soli o~wysokiej entropii atak słownikowy byłby trywialny, lecz tutaj sól tę
lukę zamyka.

\paragraph{Odporność na powtórzenie (replay).} Związanie liścia z~czwórką
\code{(registry, chainId, issuer, batchId)} uniemożliwia przeniesienie ważnego
dyplomu do innego kontraktu, sieci, wydawcy lub partii.

\paragraph{Ograniczenie.} Konstrukcja ukrywa \emph{wartości} pól, lecz nie ukrywa
\emph{struktury} koperty --- liczby pól ani faktu, że dyplom jest prywatny. Nie
zapewnia też dowodów o~ujawnionej wartości typu „stopień jest co najmniej
licencjatem'' bez podania konkretnej wartości. Takie własności wymagałyby dowodów
z~wiedzą zerową (ZK), które --- zgodnie z~przyjętym zakresem pracy --- pozostają poza
jej ramami; obecna konstrukcja realizuje selektywne ujawnianie pełnych pól, co
odpowiada założonemu wymaganiu minimalizacji danych. Funkcjonalnie odpowiada to
selektywnemu ujawnianiu znanemu z~poświadczeń weryfikowalnych
W3C~\cite{w3c-vc}, przy czym tam realizuje się je specjalizowanymi schematami
podpisów, a~tu --- wyłącznie funkcją skrótu i~solą, kosztem ujawnienia
struktury koperty.

\section{Zgodność wstecz}
\label{sec:priv-zgodnosc}

Koperta standardowa działa bez zmian --- weryfikator, który nie zna rozszerzenia
prywatnego, przetwarza ją tak jak dotąd. Na poziomie partii przyjmujemy regułę
rozłączności: dana partia jest w~całości standardowa albo w~całości prywatna, co
usuwa dwuznaczność formatu \code{docHash} przy odtwarzaniu liścia i~pozwala obu
trybom współistnieć w~jednym kontrakcie.

    \chapter{Model zarządzania uprawnieniami i analiza wartości}
\label{ch:governance}

Rozdział ten łączy dwa wątki. Najpierw domyka ostatni punkt centralizacji
w~systemie UniVerify --- pojedynczy klucz właściciela kontraktu --- zastępując go
modelem zarządzania opartym na portfelu wielopodpisowym i~mechanizmie opóźnienia
czasowego. Następnie, mając już kompletny obraz kosztów (rozdz.~\ref{ch:l2})
i~własności bezpieczeństwa systemu, przeprowadza analizę wartości: porównuje
UniVerify z~podejściem tradycyjnym oraz z~infrastrukturą klucza publicznego i~stawia
pytanie, czy przewaga blockchaina nad dojrzalszymi alternatywami uzasadnia jego
zastosowanie.

\section{Problem centralizacji uprawnień}
\label{sec:gov-problem}

W~systemie podstawowym zbiorem autoryzowanych wystawców zarządza jedno konto
zewnętrzne (\emph{externally owned account}, EOA) pełniące rolę właściciela
kontraktu \code{DiplomaRegistry}. Właściciel może dodawać i~usuwać wystawców oraz
uczelnie. Taka konstrukcja jest prosta, lecz wprowadza pojedynczy punkt awarii
o~znaczeniu krytycznym dla całego modelu zaufania.

Przejęcie tego jednego klucza pozwala napastnikowi w~jednej transakcji i~ze
skutkiem natychmiastowym przepisać zbiór uprawnionych wystawców: autoryzować
własny adres i~wydawać fałszywe, lecz formalnie poprawne dyplomy, albo odebrać
uprawnienia legalnej uczelni. Co istotne, zmiana taka jest \emph{cicha} --- dalej
w~tym rozdziale termin ten oznacza zmianę nieodróżnialną od uprawnionej operacji
administracyjnej i~niewykrywalną, zanim nabierze mocy. Osłabia to zasadniczą tezę
pracy: rejestr, który
jeden podmiot może niepostrzeżenie zmodyfikować, nie jest w~praktyce
istotnie odporniejszy na manipulację niż baza danych pod kontrolą tego samego
podmiotu.

Problem nie dotyczy przy tym pojedynczego dyplomu --- te są zakotwiczone
w~niezmiennych korzeniach Merkle --- lecz \emph{warstwy uprawnień} decydującej
o~tym, czyje korzenie są uznawane za wiarygodne. To właśnie ta warstwa wymaga
wzmocnienia.

\section{Model zarządzania: multisig + timelock}
\label{sec:gov-model}

Proponowany model przenosi własność rejestru z~pojedynczego EOA na dwuwarstwową
konstrukcję, nie zmieniając przy tym ani jednego bajtu kontraktu
\code{DiplomaRegistry}. Instalacja odbywa się wyłącznie przez istniejący,
dwustopniowy mechanizm przekazania własności
(\code{transferOwnership} $\rightarrow$ \code{acceptOwnership}), dzięki czemu
pomiary kosztów z~rozdziału~\ref{ch:l2} pozostają w~mocy.

\paragraph{Warstwa opóźnienia.} Nowym właścicielem rejestru zostaje kontrakt
\code{TimelockController} z~biblioteki OpenZeppelin~\cite{oz-contracts}. Każda
operacja administracyjna musi najpierw zostać \emph{zaplanowana}, po czym może być
\emph{wykonana} dopiero po upływie zadanego opóźnienia (\code{minDelay}). Rola
wykonawcy jest otwarta --- po upływie opóźnienia zaplanowaną operację może wywołać
dowolny podmiot --- natomiast wyłącznym \emph{proposerem} i~\emph{cancellerem} jest
portfel wielopodpisowy opisany niżej. Adres administratora timelocka ustawiono na
\code{address(0)}, co trwale usuwa administracyjną furtkę.

\paragraph{Warstwa progu.} Jedynym podmiotem uprawnionym do planowania
i~anulowania operacji jest \code{RegistryMultisig} --- portfel $m$-z-$n$ napisany
na potrzeby tej pracy, w~konwencji klasycznego \code{MultiSigWallet} firmy
ConsenSys~\cite{consensys-multisig}. Wybór implementacji własnej zamiast
audytowanego standardu branżowego (Safe) jest podyktowany celem pomiarowym:
Safe zbiera podpisy \emph{poza łańcuchem} i~wykonuje operację jedną transakcją,
przez co narzut progu staje się w~pomiarze niewidoczny. Konstrukcja użyta tutaj
zapisuje na łańcuchu każdy krok --- zgłoszenie, potwierdzenia, wykonanie ---
i~dzięki temu pozwala zmierzyć pełny, jawny koszt zarządzania
(sekcja~\ref{sec:gov-koszt}). Cena tej przejrzystości jest jednak realna: koszt
zgłoszenia operacji jest o~rząd wielkości wyższy niż w~Safe, a~implementacja
własna nie przeszła audytu i~nosi ograniczenie rezydualne opisane niżej. Do
wdrożenia produkcyjnego właściwym wyborem jest Safe
(rozdz.~\ref{ch:podsumowanie}); liczby z~sekcji~\ref{sec:gov-koszt} należy wówczas
czytać jako \emph{górne ograniczenie} narzutu. Transakcja wymaga zebrania $m$~potwierdzeń
spośród $n$~bieżących właścicieli, zanim zostanie wykonana. Portfel jest
\emph{samorządny}: operacje \code{addOwner}, \code{removeOwner}
i~\code{changeThreshold} może wywołać wyłącznie sam kontrakt (poprzez wywołanie
\code{onlySelf}), a~więc tylko za zgodą progu. Umożliwia to rotację skompromitowanego
sygnatariusza bez wymiany całej infrastruktury.

\paragraph{Przyjęte parametry.} Model jest sparametryzowany trójką
$(\code{minDelay}, m, n)$ i~bez podania jej wartości nie da się ocenić żadnej
z~gwarancji omawianych niżej. Jako konfigurację docelową przyjęto
$\code{minDelay} = 48$~godzin przy progu $m=3$ z~$n=5$. Dobór opóźnienia jest
kompromisem: musi być na tyle długie, by niezależny obserwator zdążył zauważyć
zdarzenie \code{CallScheduled} i~zaalarmować sygnatariuszy w~cyklu pracy
instytucji (dwie doby obejmują weekend), a~zarazem na tyle krótkie, by rutynowa
zmiana listy wystawców nie paraliżowała administracji. Próg $3$ z~$5$ toleruje
niedostępność dwóch sygnatariuszy bez utraty żywotności i~wymaga przejęcia trzech
niezależnych kluczy do złamania bezpieczeństwa.

Pomiary kosztu w~sekcji~\ref{sec:gov-koszt} wykonano w~konfiguracji minimalnej
$m=n=2$; narzut rośnie z~$n$, ponieważ \code{execute} przelicza liczbę
potwierdzeń, iterując po bieżącym zbiorze właścicieli. We wdrożeniu testowym na
sieci Sepolia użyto $\code{minDelay} = 120$~sekund --- wartości wyłącznie
technicznej, pozwalającej przejść pełną ścieżkę zarządzania w~jednej sesji
testowej i~nieodpowiadającej żadnemu założeniu bezpieczeństwa.

\paragraph{Kto trzyma klucze.} Model nie \emph{usuwa} zaufania, lecz je
\emph{rozprasza} --- i~uczciwość wymaga wskazania, dokąd. Ktoś musi dysponować
$n$~kluczami, a~sposób wyboru i~wymiany tego zbioru pozostaje poza łańcuchem, więc
to on wyznacza faktyczny korzeń zaufania systemu. Naturalnym kandydatem jest
konsorcjum uczelni, w~którym każda instytucja wnosi jednego sygnatariusza:
zbiór jest wtedy jawny, a~żadna pojedyncza uczelnia ani organ nadzoru nie
kontroluje progu. Rozwiązaniem alternatywnym jest powierzenie zbioru organowi
centralnemu, co jednak odtwarza pojedynczy punkt zaufania w~nowym miejscu
i~znosi większość korzyści z~całej konstrukcji. Praca nie rozstrzyga tego
wyboru --- jest on instytucjonalny, nie techniczny --- lecz wskazuje, że bez
niego twierdzenie o~,,usunięciu pojedynczego punktu zaufania'' pozostaje
niepełne: przenosi je z~jednego klucza na regulamin zbioru sygnatariuszy.

Złożenie obu warstw daje następującą własność: żadna zmiana uprawnień nie jest
ani jednoosobowa, ani natychmiastowa. Wymaga $m$~podpisów \emph{oraz} przetrwania
okna opóźnienia, w~którym pozostaje publicznie widoczna i~odwracalna
(rys.~\ref{fig:governance}).

\begin{figure}[H]
\centering
\begin{tikzpicture}[node distance=12mm and 22mm,
                    blok/.append style={minimum width=35mm},
                    blokszary/.append style={minimum width=35mm},
                    blokakcent/.append style={minimum width=35mm}]
  \node[blokakcent] (ms) {\code{RegistryMultisig}\\próg $m$-z-$n$};
  \node[blok, above=of ms] (sig) {sygnatariusze ($n$ portfeli)};
  \node[blokakcent, right=of ms] (tl) {\code{TimelockController}\\opóźnienie \code{minDelay}};
  \node[blokszary, right=of tl] (reg) {\code{DiplomaRegistry}\\(owner = timelock)};
  \draw[strz] (sig) -- node[etyk, right=1mm] {\code{submit}, \code{confirm}$\times m$} (ms);
  \draw[strz] (ms) -- node[etyk, above=2mm] {\code{schedule}} (tl);
  \draw[strz] (tl) -- node[etyk, above=2mm] {\code{execute}\\{\footnotesize po opóźnieniu}} (reg);
  \draw[strz, dashed] (ms.south) to[out=-35, in=-145]
    node[etyk, below=2mm] {\code{cancel} w~oknie opóźnienia} (tl.south);
\end{tikzpicture}
\caption[Model zarządzania: multisig + timelock]{Model zarządzania: każda operacja administracyjna wymaga progu
$m$~podpisów, a~następnie przetrwania publicznego okna opóźnienia, w~którym może
zostać anulowana. Zaplanowanie operacji emituje zdarzenie \code{CallScheduled},
dzięki czemu zmiana jest publicznie widoczna, zanim nabierze mocy.}
\label{fig:governance}
\end{figure}

\section{Model zagrożeń}
\label{sec:gov-zagrozenia}

Tabela~\ref{tab:gov-threat} zestawia model pojedynczego EOA z~proponowanym modelem
$m$-z-$n$ z~opóźnieniem. Kluczowa różnica nie polega na tym, że nowego modelu
\emph{nie da się} zaatakować, lecz na tym, że atak przestaje być cichy,
natychmiastowy i~nieodwracalny.

\begin{table}[H]
\centering
\caption{Porównanie modeli zagrożeń warstwy uprawnień.}
\label{tab:gov-threat}
\begin{tabular}{@{}lcc@{}}
\toprule
Własność & Pojedynczy EOA & Multisig 3-z-5 + timelock 48~h \\
\midrule
Wyciek jednego klucza przepisuje uprawnienia & tak & nie --- wymaga 3~podpisów \\
Zmiana uprawnień cicha i~natychmiastowa & tak & nie --- planowana i~opóźniona \\
Zmiana wykrywalna, zanim zadziała & nie & tak --- w~oknie 48~h od \code{CallScheduled} \\
Złośliwa zmiana odwracalna & nie & tak --- \code{cancel} w~tym samym oknie \\
Odtworzenie po kompromitacji sygnatariusza & --- & tak --- \code{removeOwner}/\code{addOwner} \\
Żywotność przy niedostępnym sygnatariuszu & --- & toleruje 2~niedostępnych z~5 \\
\bottomrule
\end{tabular}
\end{table}

Trzy środkowe wiersze są sednem argumentu tej pracy. Zaplanowanie operacji emituje
na łańcuchu zdarzenie \code{CallScheduled}, a~opóźnienie \code{minDelay} tworzy okno,
w~którym każdy obserwator może wykryć nadchodzącą zmianę, zanim ta nabierze mocy.
W~tym samym oknie próg sygnatariuszy może operację anulować. Rejestr uprawnień staje
się więc \emph{publicznym, dopisywalnym i~odpornym na ciche wsteczne przepisanie}
logiem --- dokładnie tą własnością, do której wraca analiza wartości
w~sekcji~\ref{sec:gov-porownanie}.

\paragraph{Cena opóźnienia: reakcja awaryjna.} Opóźnienie chroni przed cichą
zmianą uprawnień, lecz w~jednym scenariuszu działa przeciw obrońcy. Gdy dojdzie
do przejęcia klucza \emph{wystawcy}, odebranie mu uprawnień wymaga zebrania
$m$~podpisów \emph{oraz} przeczekania \code{minDelay} --- przy przyjętych 48~h
napastnik dysponuje dwiema dobami na publikowanie kolejnych fałszywych korzeni.
Sytuację pogarsza brak unieważniania na poziomie partii: kontrakt udostępnia
wyłącznie \code{revokeFromBatch}, działające na pojedynczym dokumencie, więc
wycofanie sfałszowanej partii 10~tys.\ dyplomów wymagałoby 10~tys.\ transakcji.

Jest to realna luka projektowa, nie tylko brak opisu. Wskazać można dwie ścieżki
jej domknięcia, obie pozostawione poza zakresem implementacyjnym pracy. Pierwsza
to \code{revokeBatchRoot} --- operacja unieważniająca cały korzeń jednym zapisem,
tania i~prosta, lecz gruboziarnista: cofa również dyplomy wydane w~tej partii
uczciwie. Druga to wydzielenie roli \emph{strażnika} (\emph{guardian})
z~uprawnieniem wyłącznie do anulowania zaplanowanych operacji i~wstrzymania
wystawcy, bez prawa zmiany rejestru --- ścieżka szybka, bo pomijająca opóźnienie,
a~zarazem niegroźna, bo z~definicji nie może niczego dodać. Druga jest właściwsza:
zachowuje asymetrię, w~której działania \emph{rozszerzające} uprawnienia są powolne
i~jawne, a~\emph{ograniczające} --- natychmiastowe.

\paragraph{Ograniczenie rezydualne.} Operacja \code{removeOwner} nie czyści
potwierdzeń usuniętego właściciela na transakcjach nadal oczekujących. Jeśli ten sam
adres zostanie później ponownie dodany, jego stare potwierdzenie zaczyna liczyć się
na nowo. Jest to standardowy kompromis wzorca ConsenSys, unikający nieograniczonej
iteracji po zbiorze oczekujących transakcji przy każdym usunięciu. Operacyjnie
należy unikać ponownego dodawania byłego właściciela, gdy istnieją wrażliwe
transakcje oczekujące.

\section{Koszt zarządzania}
\label{sec:gov-koszt}

Narzut modelu zmierzono testem \code{GovernanceGasTest}
(\code{forge test --gas-report}) w~konfiguracji minimalnej $m=n=2$; dla progu
docelowego 3-z-5 narzut jest wyższy o~koszt iteracji po większym zbiorze
właścicieli w~\code{execute}. Wszystkie wdrożenia kontraktów oraz jednorazowe
przekazanie własności realizowane są w~\code{setUp()}, dzięki czemu ciało każdego
testu mierzy dokładnie jedną operację onboardu wystawcy, a~raportowana przez Foundry
suma na test (wyłączająca \code{setUp}) jest liczbą podaną bezpośrednio, niewymagającą
dekompozycji.

Linią bazową sprzed wprowadzenia zarządzania jest pojedyncze wywołanie
\code{onboardIssuer\allowbreak AndUniversity} na rejestrze będącym własnością zwykłego EOA:
84\,246 gazu. Ta sama operacja przeprowadzona pełną ścieżką zarządzaną ---
\code{multisig.submit} + \code{multisig.confirm} + \code{multisig.execute}
(planuje na timelocku) oraz \code{timelock.execute} (wywołuje rejestr po
opóźnieniu) --- kosztuje 824\,254 gazu, co daje narzut na zmianę uprawnień rzędu
$824\,254 - 84\,246 = 740\,008$ gazu.

Tabela~\ref{tab:gov-cost} pokazuje, gdzie trafia gaz ścieżki zarządzanej. Dominuje
\code{submit}, który zapisuje na łańcuchu pełne dane zaplanowanej operacji.

\begin{table}[H]
\centering
\caption{Rozbicie kosztu zarządzanego onboardu (wykonanie on-chain).}
\label{tab:gov-cost}
\begin{tabular}{@{}lrl@{}}
\toprule
Wywołanie & Gaz & Uwaga \\
\midrule
\code{RegistryMultisig.submit}   & 487\,647 & dominuje --- zapis danych operacji \\
\code{RegistryMultisig.confirm}  &  52\,213 & niezależne od danych \\
\code{RegistryMultisig.execute}  & 160\,119 & wywołuje \code{timelock.schedule} \\
\code{TimelockController.execute} & 99\,240 & wywołuje rejestr \\
\midrule
Suma & 799\,219 & samo wykonanie on-chain \\
\bottomrule
\end{tabular}
\end{table}

Suma czterech wywołań (799\,219 gazu) jest niższa od kosztu całego testu
(824\,254 gazu) o~narzut ciała testu --- przygotowanie i~kodowanie ABI
argumentów wywołań, około 3\% --- który nie jest kosztem on-chain. Porównanie
z~linią bazową używa spójnie sum na~test po~obu stronach, więc narzut ten
znosi się w~różnicy jedynie częściowo i~podana wartość 740\,008 gazu jest
oszacowaniem ostrożnym (z~góry).

Istotne jest przy tym, że zmiany uprawnień (dodanie lub usunięcie wystawcy) są
\emph{rzadkie}. Narzut 740\,008 gazu ponoszony jest raz na wystawcę i~amortyzuje się
na wszystkich dyplomach, które ten wystawca później wyda. Koszt rejestracji
pojedynczego dyplomu, zmierzony w~rozdziale~\ref{ch:l2}, pozostaje więc
niezmieniony.

\section{Porównanie podejść: status quo, PKI, blockchain}
\label{sec:gov-porownanie}

Mając kompletny obraz kosztów i~gwarancji, można postawić pytanie zasadnicze: co
właściwie daje blockchain wobec tańszych i~dobrze zrozumianych alternatyw?
Poniżej zestawiono trzy podejścia do weryfikacji dyplomu:

\begin{description}
  \item[Status quo] --- weryfikator kontaktuje się z~uczelnią (telefon, e-mail,
        formularz) i~prosi o~potwierdzenie. Dominujący dziś sposób.
  \item[PKI] --- uczelnia podpisuje cyfrowo dokument dyplomu kluczem
        certyfikowanym w~infrastrukturze klucza publicznego; weryfikator sprawdza
        podpis offline wobec certyfikatu wystawcy.
  \item[PKI + CT] --- to samo, uzupełnione o~publiczny, wyłącznie dopisywalny
        log wystawionych certyfikatów wystawców w~modelu \emph{Certificate
        Transparency}~\cite{rfc6962} (sekcja~\ref{sec:pki}).
  \item[Blockchain (UniVerify)] --- korzenie Merkle partii oraz zbiór uprawnionych
        wystawców publikowane są w~kontrakcie na publicznym blockchainie.
\end{description}

Tabela~\ref{tab:gov-value} zestawia je względem sześciu własności.

\begin{table}[H]
\centering
\caption{Macierz wartości: cztery podejścia wobec sześciu własności.}
\label{tab:gov-value}
\begin{tabular}{@{}lcccc@{}}
\toprule
Własność & \makecell{Status\\quo} & PKI & \makecell{PKI\\+ CT} & UniVerify \\
\midrule
Weryfikacja bez kontaktu z~wystawcą                & nie & tak & tak & tak \\
Niski koszt jednostkowej weryfikacji              & nie & tak & tak & tak \\
Działa, gdy wystawca przestał istnieć              & nie & częściowo & częściowo & tak \\
Publiczny, dopisywalny log uprawnień do~wydawania & nie & nie & tak & tak \\
Odporność na ciche wsteczne odebranie uprawnień   & --- & nie & tak & tak \\
Jeden widok logu bez warstwy dodatkowej           & --- & --- & nie & tak \\
\bottomrule
\end{tabular}
\end{table}

Ocena „częściowo'' w~trzecim wierszu wymaga komentarza. Podpis cyfrowy
pozostaje weryfikowalny offline tak długo, jak długo daje się zwalidować
łańcuch certyfikatów: po zniknięciu wystawcy certyfikaty wygasają bez
możliwości odnowienia, a~usługi statusu (CRL/OCSP) przestają być utrzymywane.
Weryfikacja długoterminowa wymaga wówczas dodatkowej infrastruktury
znakowania czasem i~archiwizacji materiału walidacyjnego (podpisy w~formatach
długoterminowych), która w~praktyce rzadko obejmuje dyplomy --- stąd ocena
częściowa, nie negatywna.

Zasadniczy wniosek z~tabeli jest dla blockchaina niewygodny, lecz uczciwy. PKI
\emph{dorównuje} rozwiązaniu blockchainowemu w~dwóch pierwszych, najczęściej
przywoływanych zaletach: pozwala na niezależną weryfikację offline bez kontaktu
z~wystawcą i~po znikomym koszcie jednostkowym. Podpis cyfrowy pod dyplomem daje to
samo bez konieczności utrzymywania blockchaina. Argument „blockchain eliminuje
telefon do dziekanatu'' nie jest więc argumentem za \emph{blockchainem} --- jest
argumentem za \emph{kryptografią asymetryczną}, którą PKI dostarcza taniej.

\section{Czy blockchain tu wygrywa?}
\label{sec:gov-wniosek}

Różnica pojawia się dopiero w~dwóch ostatnich wierszach tabeli~\ref{tab:gov-value},
dotyczących nie pojedynczego dokumentu, lecz \emph{warstwy uprawnień do wydawania}.
W~modelu PKI o~tym, które certyfikaty wystawców są ważne, decyduje urząd certyfikacji.
Może on unieważnić lub wystawić certyfikat prywatnie, a~listy unieważnień (CRL/OCSP)
odzwierciedlają \emph{bieżący} stan --- nie dają publicznego, niemodyfikowalnego
zapisu tego, kto i~kiedy był uprawniony, ani nie chronią przed cichą wsteczną
zmianą tego zbioru. Weryfikator ufa, że urząd nie przepisał historii.

\paragraph{Certificate Transparency zamyka tę lukę --- prawie.} Powyższy argument
byłby rozstrzygający, gdyby publiczny log uprawnień dało się uzyskać wyłącznie
na blockchainie. Tak nie jest. Certificate Transparency~\cite{rfc6962} realizuje
dokładnie tę własność --- publiczny, wyłącznie dopisywalny log wystawionych
certyfikatów z~dowodami spójności --- na drzewach Merkle, bez łańcucha bloków,
produkcyjnie od ponad dekady i~po koszcie nieporównanie niższym niż utrzymanie
kontraktu. Czwarta i~piąta własność w~tabeli~\ref{tab:gov-value} przestają zatem
odróżniać blockchain od dojrzałej alternatywy. Uczciwe postawienie sprawy wymaga
przyznania, że \emph{większość} przewagi, którą praca zamierzała wykazać, jest
osiągalna w~modelu PKI+CT.

Pozostaje różnica wąska i~jedna, ujęta w~ostatnim wierszu tabeli. Log CT jest
prowadzony przez wyznaczony zbiór operatorów, a~o~tym, którym z~nich klient ufa,
rozstrzyga polityka dostawcy przeglądarki. Log może wobec różnych obserwatorów
przedstawić \emph{rozdwojony widok} (\emph{split view}): dwa niesprzeczne
wewnętrznie, lecz różne stany, z~których każdy ma poprawne dowody spójności.
Wykrycie takiego rozdwojenia wymaga, by obserwatorzy porównywali swoje widoki
między sobą. Mechanizm do tego przeznaczony był specyfikowany w~IETF jako
osobny dokument~\cite{trans-gossip}, lecz nigdy nie osiągnął statusu standardu
--- prace nad nim wygasły. W~praktyce spójność widoku opiera się na wymogu
zebrania poświadczeń od kilku niezależnych operatorów oraz na ich wzajemnym
audycie, a~więc na założeniu, że operatorzy nie działają w~zmowie.

Publiczny łańcuch bloków rozwiązuje ten sam problem inaczej: jednolitość widoku
jest u~niego własnością warstwy konsensusu, nie dodatkową warstwą protokołu.
Nie ma logu, którego stan mógłby się rozdwoić wobec dwóch weryfikatorów, bo
o~stanie rozstrzyga ten sam mechanizm, który utrzymuje sam łańcuch.

Marginalna wartość blockchaina w~tym zastosowaniu sprowadza się zatem do
własności węższej, niż wynikałoby z~porównania z~samym PKI: \emph{publicznego,
dopisywalnego, odpornego na manipulację logu tego, kto może wydawać dyplomy,
o~jednym widoku dla wszystkich weryfikatorów, bez zaufania do zbioru operatorów
logu}. Logu, którego nie da się
po cichu wstecznie przepisać. Model zarządzania z~sekcji~\ref{sec:gov-model} jest
tym, co tę własność urzeczywistnia --- bez timelocka i~progu sam blockchain tej
gwarancji by nie dawał, bo pojedynczy właściciel mógłby log przepisać w~granicach
swoich uprawnień.

\paragraph{Studium przypadku: uczelnia skali politechniki.} Skalę kosztu dobrze
obrazuje rachunek dla dużej uczelni technicznej wydającej rzędu 10~tys.\ dyplomów
rocznie. Publikacja partii na~Base kosztuje około $9\cdot10^{-4}$~USD niezależnie
od jej rozmiaru (tab.~\ref{tab:issue-tx}). Nawet przy kilkunastu partiach rocznie
--- na przykład osobnej partii na wydział i~semestr --- łączny koszt publikacji nie
przekracza dwóch centów. Po doliczeniu jednostkowych unieważnień
($1{,}1\cdot10^{-3}$~USD, tab.~\ref{tab:revoke}) oraz sporadycznych,
amortyzowanych zmian uprawnień (sekcja~\ref{sec:gov-koszt}) całkowity roczny koszt
on-chain pozostaje poniżej dziesięciu centów.

Koszt transakcyjny nie jest zatem czynnikiem decyzyjnym w~żadnym realistycznym
scenariuszu wdrożenia. Realne koszty leżą w~integracji z~systemami dziekanatowymi
i~w~utrzymaniu operacyjnym (rozdz.~\ref{ch:podsumowanie}), a~wybór między
rozważanymi podejściami rozstrzyga się na osi wymagań, nie kosztów.

Pozostaje pytanie o~proporcję. Wyodrębniona wyżej własność kosztuje --- zależnie
od sieci --- ułamki mikrocenta na dyplom (rozdz.~\ref{ch:l2}) plus jednorazowy
narzut 740\,008 gazu na zmianę uprawnień. Jest to koszt niski w~wartościach
bezwzględnych, lecz niezerowy --- i~ponoszony wobec dwóch alternatyw, z~których
każda jest dojrzalsza.

Praca nie rozstrzyga tej proporcji arbitralnie. Wniosek ma postać stopniowanej
rekomendacji, a~nie opowiedzenia się po jednej stronie.

\begin{description}
  \item[Samo PKI wystarcza,] gdy wymaganiem jest wyłącznie niezależna od wystawcy,
        tania weryfikacja dokumentu. Argument ,,blockchain eliminuje telefon do
        dziekanatu'' jest w~całości zaspokojony na tym poziomie.
  \item[PKI z~Certificate Transparency wystarcza,] gdy dodatkowym wymaganiem jest
        publiczna, historyczna audytowalność uprawnień do wydawania --- i~jest to
        rozwiązanie właściwe dla zdecydowanej większości realnych wdrożeń,
        tańsze i~lepiej zbadane niż blockchain.
  \item[Blockchain wygrywa wtedy i~tylko wtedy,] gdy do powyższego dochodzi
        wymaganie \emph{jednolitości widoku bez zaufania do operatorów logu}:
        gdy weryfikatorzy nie mogą przyjąć, że wyznaczony zbiór operatorów nie
        działa w~zmowie i~nie przedstawia im rozdwojonego widoku, ani nie mogą
        polegać na warstwie uzgadniania widoków, która w~ekosystemie CT nie
        została ustandaryzowana. Scenariuszem modelowym jest weryfikacja
        ponadinstytucjonalna i~transgraniczna, w~której żaden pojedynczy podmiot
        --- ani urząd certyfikacji, ani operator logu --- nie jest akceptowany
        przez wszystkie strony.
\end{description}

Trzy powyższe progi, a~zwłaszcza wyodrębnienie trzeciego, stanowią główny wkład
analityczny niniejszej pracy. Jest on dla tezy blockchainowej mniej korzystny,
niż zakładano na wstępie: uwzględnienie Certificate Transparency przenosi
większość deklarowanej przewagi na rozwiązanie bez łańcucha bloków i~zawęża
uzasadnione zastosowanie blockchaina do przypadku brzegowego. Przypadek ten
pozostaje jednak realny, a~jego cena transakcyjna --- znikoma w~porównaniu
z~kosztem sporu o~to, komu wolno wystawiać dyplomy.

    \chapter{Podsumowanie i dalsze prace}
\label{ch:podsumowanie}

\section{Wnioski}
\label{sec:wnioski}
% Podsumuj trzy wkłady i główny wynik (odpowiedź na pytanie "czy warto").
\TODO{Zwięzłe wnioski końcowe spinające rozdziały 4--7.}

\section{Ograniczenia}
\label{sec:ograniczenia}
% Zbierz ograniczenia z rozdziałów: pomiary L2, małe domeny w prywatności,
% pojedynczy region pomiaru, założenia modelu kosztowego.
\TODO{Uczciwa lista ograniczeń pracy.}

\section{Dalsze prace}
\label{sec:dalsze-prace}
% Odrzucone świadomie w tej pracy (poza zakresem): dowody ZK (predykatowe,
% "GPA > 3.0"), akumulatory / drzewa Verkle, weryfikacja cross-chain,
% produkcyjne zarządzanie kluczami studenta.
\TODO{Kierunki rozszerzeń.}


    % Bibliografia (plik EE-dyplom.bib, biblatex + biber)
    \bibliografia

    % Strony końcowe
    \acronymslist
    \listoffigures
    \listoftables
    % \easyappendices        % wymaga tekst/appendices.tex
\end{document}
